\section{Research log}

\subsection*{5 September: smooth literal fibres and exact period/cusp refinements}
Section~\ref{sgd:smooth-addon:section} proves that the literal
scalar--Gauss domain is embedded with exactly four six-dimensional
components and gives
$\mathcal S_{\rm sg}\cong\{1,-1\}^2\times(0,\infty)
\times\mathbb R^2\times\mathbb T^3$.
The first continuous coordinate is the unchanged $t=g^2$; each specified
positive level has four embedded five-dimensional components
$\mathbb R^2\times\mathbb T^3$. The inverse uses
$v_0=\sqrt5\sqrt{\delta^2+y^2}/2+e^s$ and
$u=(4\tau v_0+72(\delta^2+y^2))/(72v_0+5\tau)$,
with $\delta=\sqrt{2/3}$ and $\tau=t$.
The full literal inverse and the phase map
$(\zeta,p_+,p_-)\mapsto(\zeta^2,p_+/\zeta,p_-/\zeta)$ are proved
smooth, retaining both representatives of the sole parameter kernel.
No spatial equation or physical gauge quotient is imposed, and no quantum
gap follows. The shared source is included once after the existing
scalar--Gauss source in the common quaternionic module.

Proposition~\ref{prop:period-exact-admissible-constants} proves the exact
half-plane $\operatorname{Im}c_0<-M$, where
$M=\max_{B\setminus\{p_0\}}D_{\rm part}$ is finite and attained.
The cusp-free affine solution $\beta+i\,t\circ\pi$ has
$\widetilde D=D+\operatorname{Re}t$ unbounded above and is excluded by
the actual cusp extension. Equation~\eqref{eq:cusp-exact-opposite-edge}
computes $\iota_v=\Psi_{B_0v}|_{C_v}$ with multipliers
$e^{-2\pi i\mu_0}$, $e^{12\pi i\mu_0}$ and
$e^{-2\pi i(7\mu_0+b_0)}$ in the stated same-character conventions.
The cusp proof also restores restricted chart domains, a bound on a fixed
closed subdisc, connectedness of $Y_\epsilon$, and weak inequalities for
closed polydiscs. The bounded finite-filling and gluing recheck preserves
the actual logarithmic signs and complete overlap maps. These audits
add no sphere-identification or CDP theorem. Historical annotations remain
unchanged; builds and publication have separate receipts.

\subsection*{5 September 2026 --- higher wrap, octonionic carrier, and selector bridge}
The complete draft in Section~\ref{sec:h24-continuation} retains the
higher cusp lattice quotient and its local torus morphisms, constructs
the octonionic determinant-boundary pair $(B^{25},S^{24})$, and writes
the comparison with $F_4/\operatorname{Spin}(8)$ and its six fixed points.
The peripheral character factors through $SL_2(\mathbb Z)_{\mathrm{ab}}
\cong\mathbb Z/12$. Sixfold suspension identifies this order-twelve
homotopy group with the even subgroup in the stable order-twenty-four
group. The particular cusp coefficient is uncomputed. The source audit
corrected the flag pair/orbit tangent identification, root factors of two,
scope of the finite $S_3$ comparison, and the order of the Euler class
after coefficient reduction. The reduction has order $N/\gcd(N,6)$;
no classification of all possible arithmetic invariants is asserted.
The $p=1201$ section proves the exact three-dimensional Jordan span and
preserves the independent full census and same-residue comparison counts.

\subsection*{Round 0 --- workbench initialization}

\begin{itemize}[leftmargin=1.8em]
\item Preserved the completed 121-page annotated audit as
\Source{supporting_materials/workbench_history/former_workbench_checkpoint_56_pages/audit_baseline/annotated_audit.tex}
and the 108-page reconstruction as
\Source{supporting_materials/workbench_history/former_workbench_checkpoint_56_pages/audit_baseline/clean_transcription.tex}.
\item Converted the review objections into named mathematical obligations with explicit maps, bases, signs, and falsification tests.
\item Proved the exact conditional endpoint: once the geometric maps are identified with the printed matrices, the verified algebra gives the standard smooth $6$-sphere.
\item Kept the conductor mechanism separate from both the existence of $X$ and the truth or falsity of the published CDP theorem.
\end{itemize}

\subsection*{Round 1 --- validate, repair, or disprove reset}

\begin{itemize}[leftmargin=1.8em]
\item Removed the privacy label and made the notebook suitable for later sharing without changing its present Overleaf permissions.
\item Reclassified the actual assertion $X\cong_{\rm diff}S^6$ as \Open. Reserved \Conditional\ for the separately proved implication from the exact map-identification obligations C1--C3, V1, MV1, and L1.
\item Made every ChatGPT/Codex objection falsifiable: a valid outcome is that the manuscript's named proof discharges an obligation and the reviewer objection fails.
\item Adopted a literature gate: each literature-bearing calculation begins by inventorying and reading the relevant local corpus, recording exact hypotheses and locators before new deductions are attempted.
\item Compared the earlier Fable/JGPTech Jacobian certificate and Erd\H{o}s--Strauss Jacobi-symbol work with the present problem. No theorem or typed map transfers to $S^6$; the exact transfer is the marked-data/boundary-closure/typed-map audit architecture recorded in
\Source{supporting_materials/workbench/comparative_context/prior_fable_context.tex}.
\end{itemize}

\subsection*{Current next calculations}

\begin{enumerate}[leftmargin=1.8em]
\item Obtain an independent specialist replication of the now-complete
cotangent direct-image calculation.  The exact marked-fibre
period-coordinate representatives, global Bott--Chern generators, and all
their products are now complete and are no longer active obligations.
\item The third proof that
\(\eta^*i^*:\operatorname{Pic}(X)\to\operatorname{Pic}(\mathrm{dP}_6)\)
vanishes is now complete.  It uses $\mathcal O^*$ cocycles, both finite-filling
norm identities, the cusp argument principle, and toric support functions,
while avoiding both the global Leray differential and local integral
specialization/bounding-chain input.  Continue by testing the next
candidate-specific direct-image-length and conductor-localization questions.
\item The corrected-replacement problem for CDP20 Theorem~2.2 is now closed
by \cref{thm:corrected-cdp-global,thm:corrected-cdp-fibrewise}.  Continue by
testing whether the positive Euler characteristic admits a quantitative
localization in the torsion length of \(R^2f_*(T_X\otimes L)\) and the
conductor connecting kernels, and invite independent specialist replication
of the candidate counterexample.  The global Mayer--Vietoris calculation is
also closed and is no longer an active obligation.
\end{enumerate}

\subsection*{Round 2 --- local-primary-source audit}

\begin{itemize}[leftmargin=1.8em]
\item Read the local primary SGA~7~I and SGA~7~II material at the exact
semistable-reduction, monodromy, nearby-cycle, specialization, and proper
direct-image locators recorded in
\Source{supporting_materials/workbench/research/literature_log.tex}.
\item Confirmed that this formalism does not, by citation alone, prove the
manuscript's integral injectivity and equality-with-invariants assertion.
\item Located smooth complex-manifold Serre duality and algebraic
cohomology-and-base-change results, and recorded precisely why neither
certifies the singular analytic fibre step as presently written.
\item Verified the final smooth six-sphere recognition route conditional on
the construction, simple-connectivity, and integral-homology inputs.
\item Flagged the absent CDP and analytic-duality/base-change primaries as
explicit source gates rather than filling them with recollection.
\end{itemize}

\subsection*{Round 3 --- C1 objection corrected and reduced}

\begin{itemize}[leftmargin=1.8em]
\item Re-read Proposition~7.2 and its proof in Lemmas~7.5--7.10.  The
manuscript does provide a detailed equivariant-collapse construction; the
review formulation ``construct a finite CW model'' was too strong and is not
a logically necessary replacement for that proof.
\item Proved independently the formal topological spine of Lemmas~7.8--7.9:
under the printed contractibility, free-action, compact-quotient, common-deck-
group, and collar hypotheses, the boundary quotient is a strong deformation
retract and the retraction lifts equivariantly.
\item Narrowed C1 to the candidate-specific facts that still require audit:
Theorem~4.5, continuity of the straightening map at the toric boundary,
manifold and compactness assertions for the positive quotient, stabilizers,
and the phase-spreading construction.  C2--C3, V1, MV1, and L1 now cite their
actual manuscript proof chains instead of being described as absent.
\end{itemize}

\subsection*{Round 4 --- CDP primary comparison and analytic bridges}

\begin{itemize}[leftmargin=1.8em]
\item Acquired and read the official CDP20 journal article, the
author-hosted CDP98 article, and the Ramis--Ruget duality paper; inspected
B\u{a}nic\u{a}--St\u{a}n\u{a}\c{s}il\u{a}, Chapter~III, Section~3, at the
exact analytic base-change theorem and corollary.  All copies remain outside
the shareable project and exact paths, locators, URLs and hashes are in
\Source{supporting_materials/workbench/research/ledgers/literature_sources.md}.
\item Verified that CDP20 Theorem~2.2 has exactly the hypotheses claimed for
the candidate, and isolated the p.~692 failure: the conormal sequence controls
sections mapping to zero in $\Omega^1_S$, whereas passing to
$\widetilde\Omega^1_S$ also kills nonzero conductor-supported torsion.
\item Constructed an independent proper test family from a smoothing of a
nodal plane cubic crossed with $\mathbb P^1$.  A nontrivial gluing line bundle
has trivial pullback to the normalization and gives both vanishings used on
p.~692, while the conductor differential gives a nonzero ambient section.
This refutes the printed inference without claiming a counterexample to the
global CDP theorem.
\item Closed the two formal analytic bridges left open by the earlier audit.
Cartier exact sequences and natural Serre duality on the smooth ambient
threefold prove the required duality on the nonnormal fibre; Ramis--Ruget
gives an independent singular-space check.  Top-degree analytic base change
is an isomorphism by B\u{a}nic\u{a}--St\u{a}n\u{a}\c{s}il\u{a}, Theorem~3.4
and Corollary~3.5, because fibrewise $H^3$ vanishes.
\item Consequently, conditional on the upstream construction of the smooth
compact $X$, the asserted fibre $W\simeq\mathrm{dP}_6/(\text{opposite
sides})$, and $H^2(X;\mathbb Z)=0$, the manuscript's conclusion
$R^2f_*(T_X\otimes L)\ne0$ for every $L$ is now validated.  The unresolved
question is global: a failure in the candidate, an independent repair of
CDP20, or a genuine counterexample to its theorem.
\end{itemize}

\subsection*{Round 5 --- provisional candidate proof reconstruction
(status corrected in Round 6)}

\begin{itemize}[leftmargin=1.8em]
\item Reconstructed the modular lift, the two affine torsor problems, the
determinant inequality, and the equivariant period-family formula.  A later
line audit found that local torsor nonemptiness and the full integral
equivariance calculation still had to be restored; the period family
therefore remained \Open.
\item Expanded the toric cusp action through its quantitative $B_t$
estimate, action law, bounded-chart control, proper discontinuity, and local
normal-crossings equations.  The quotient geometry is \Conditional\ on the
period input and the remaining global quotient identifications.
\item Reconstructed the finite logarithmic fillings and global gluing, but a
later audit found undefined linearizing data and omitted branch-independence
checks.  Thus (H1)--(H3) were not discharged in this round.
\item Reconstructed Appendix~B's integral nearby-cycle calculation,
including the normalization and vanishing-cycle resolutions, the
hypercohomology ranks, and the supported Gysin strategy.  Applicability of
integral analytic nearby-cycle comparison to the proper holomorphic family
and the full supported-specialization diagram remained \Open.
\item Computed the abstract global Leray page and the arithmetic of the three
candidate transgressions.  Their identification with the differentials of
the actual geometric map remained the named L1 obligation; consequently the
vanishing of $H^1,H^2,H^3$ was only \Conditional.
\item Reconstructed the abstract based van Kampen presentation from the
affine deck generators and cusp relations.  Based inclusions, primitivity,
signs and drift remained the named V1 hypotheses, so simple connectivity and
sphere recognition were only \Conditional.
\item Reconstructed the very-general fibre calculation and candidate
Lemma~9.6's meromorphic descent, adding the omitted connectivity proof and
root estimate.  The global conclusion
$\mathcal M(X)=f^*\mathbb C(t)$ and $a(X)=1$ remained
\Conditional\ at this historical round because the family had not yet been
closed.
\item Refined the CDP20 p.~692 defect to the two exact sheaf sequences
containing $\mathscr K_S$, and traced the separate semidirect-product
coinvariant error from CDP98 into the CDP20 torus branch.
\item Reclassified the cellular collapse/Mayer--Vietoris calculation as an
open alternate certificate rather than a premise of the proposed
nearby-cycle/Leray route.  Neither route was complete at this stage.
\end{itemize}

\subsection*{Round 6 --- 25 August 2026: fail-closed recalibration and
analytic-source audit}

\begin{itemize}[leftmargin=1.8em]
\item Reran the two retained exact-arithmetic harnesses: the
full-manuscript suite passed all $28$ named checks and the independently
implemented core suite passed all $38$ named checks.  These certify their
displayed algebra, not the geometric origin of the matrices.
\item Re-read the local SGA~7 corpus at the relevant nearby-cycle loci and
separated its algebraic scope from the present possibly nonalgebraic proper
holomorphic map.  Added Dimca, \S4.2, and Kashiwara--Schapira, \S8.6, as the
analytic references whose precise integral hypotheses must be checked.
\item Restored the full supported-specialization diagram: equivariant local
product neighbourhood, relative-cohomology excision, saturated subtorus
$\Lambda_{\rm tor}$, Thom isomorphism, and primitive K\"unneth Gysin
summand.  The internal deduction is complete conditional on the analytic
comparison theorem and the candidate local models.
\item Derived the global skyscraper exact sequence stalk by stalk and
expanded the parabolic-cocycle quotient.  Reclassified the global Leray page,
the three transgressions, low-degree cohomology, van Kampen presentation,
and sphere recognition as \Conditional\ with their analytic and geometric
hypotheses named explicitly.
\item Corrected the theorem-level CDP disposition.  The p.~692 implication
and the unrestricted CDP98 direct-sum formula are \Refuted; this proves two
printed proof-route defects, not by itself the falsity of CDP20
Theorem~2.2.  A theorem-level counterexample remains \Conditional\ on the
candidate construction, topology, and algebraic-reduction chains.
\item Added the exact toric fundamental-group source
Cox--Little--Schenck, Theorem~12.1.10, with Danilov, \S9, as an independent
comparison; replaced the circular exhaustion argument by a finite-subfan
null-homotopy followed by explicit $t$-coordinate scaling.
\end{itemize}

\subsection*{Round 7 --- 26 August 2026: central non-collapse proof closed}

\begin{itemize}[leftmargin=1.8em]
\item Completed the period-and-torsor descent, the periodic toric cusp
quotient, the punctured cusp overlap, both finite logarithmic fillings, the
finite radial-stalk maps, and the global Hausdorff compact gluing.  The
construction theorem is now \Verified.
\item Completed the analytic integral nearby-cycle proof at the cusp:
local Milnor fibres, normalization resolution, supported Thom--Gysin class,
saturation, and the exact invariant lattices are all explicit.  The
specialization
\[
H^q(W;\mathbb Z)\xrightarrow{\sim}H^q(F;\mathbb Z)^{T_0}
\]
is \Verified\ and is not defined through the separate cellular collapse.
\item Completed the global parabolic-sheaf calculation and identified the
Cech discrepancy with Leray \(d_2\) through the canonical truncation
triangle for \(Rf_*\mathbb Z_X\).  Multiplicativity gives all three primitive
transgressions; consequently
\[
H^1(X;\mathbb Z)=H^2(X;\mathbb Z)=H^3(X;\mathbb Z)=0,
\qquad H_2(X;\mathbb Z)=0.
\]
\item Closed the based van Kampen maps from the actual coordinate overlaps.
Zero-section whiskers, primitive clockwise meridians, finite power relations,
the cusp quotient \(\Lambda\to\bar\Lambda\), and the toric bounding disc
produce the exact terminal presentation with \(p=-1\), hence
\(\pi_1(X)=1\).  Sphere recognition and
\(\langle c_3(T_X),[X]\rangle=2\) are now \Verified.
\item Closed the very-general N\'eron--Severi/Appell--Humbert calculation
and meromorphic descent.  The constructed map satisfies
\[
\mathcal M(X)=f^*\mathbb C(t),\qquad a(X)=1,
\]
and is the algebraic reduction.
\item The preceding verified chains show that the constructed threefold
satisfies the stated hypotheses of CDP20 Theorem~2.2 while violating its
conclusions \(c_3(X)\le0\) and \(a(X)=0\).  The theorem-level counterexample
status now follows from the reconstructed mathematics; it is not inferred
only from the two separately located defects in the printed CDP route.
\item Kept C1--C3, MV1, and the Mayer--Vietoris comparison part of L1
\Open.  They constitute an alternate topological certificate and are not
dependencies of the verified intrinsic nearby-cycle/Leray route.
\item Kept the candidate's Proposition~10.8 application \Conditional.
Its period/countability strategy still requires workbench proofs of the
Section~9 direct-image and degree-one base-change statements,
\(\operatorname{Pic}(X)\simeq\mathbb C\), and the exact finite-fibre
normal-bundle orders.
\item Reviewer~B independently examined the reconstruction existing up to
that report's evidentiary cutoff.  Its provisionally favourable disposition
agrees that the non-collapse nearby-cycle/Leray and based van Kampen route
does not depend on the open collapse comparison.  The later Codex proofs
recorded in this round are separate and are not attributed to that review.
\end{itemize}

\subsection*{Round 8 --- 26 August 2026: Proposition 10.8 closed}

\begin{itemize}[leftmargin=1.8em]
\item Reconstructed the Hodge line
\(\mathfrak K=f_*\omega_{X/B}\simeq\mathcal O_{\mathbb P^1}(1)\)
from the explicit determinant characters, finite ramification, cusp toric
volume form, and modular section.
\item Proved the analytic normalization resolution of the cusp fibre and
obtained
\[
h^q(W,\mathcal O_W)=(1,2,1),\qquad
H^1(W;\mathbb C)\simeq H^1(W,\mathcal O_W),
\]
with nonzero degree-one cup product.  Intrinsic nearby-cycle
specialization replaces both uses of the candidate's open collapse route.
\item Proved local freeness and base change for every
\(R^qf_*\mathcal O_X\), including the nonreduced fibres, and reconstructed
\[
R^1f_*\mathcal O_X\simeq\mathcal O\oplus\mathcal O(-1),\qquad
R^2f_*\mathcal O_X\simeq\mathcal O(-1).
\]
The nowhere-zero global section is the Dolbeault image of the invariant
class \(\gamma\).
\item The exponential sequence and verified integral topology now give
\(\operatorname{Pic}(X)\simeq\mathbb C\).  A direct finite-cover character
argument gives
\(\operatorname{ord}(\mathcal O_X(S_j)|_{S_j})=m_j\).
\item Checked the full period calculation.  For
\(\lambda=a\gamma+bu+cw+d\delta\), membership in the Picard restriction
line is exactly
\[
c\tau+d\mu-b=0.
\]
Linear independence of \(1,\tau,\mu\) proves that all parameters outside
\(\mathbb R\) and a countable set are simultaneously non-torsion on every
scheme-theoretic fibre and on both \(S_j\).  Proposition~10.8 is now
\Verified.
\item Independent source and literature audits agreed with this promotion.
Reviewer~B's earlier report is retained as an independent review only up
to its own evidentiary cutoff; the proof in this round is subsequent work.
\end{itemize}

\subsection*{Round 9 --- 26 August 2026: public synchronization}

\begin{itemize}[leftmargin=1.8em]
\item Rebuilt the workbench after the Proposition~10.8 promotion.  The local
  \(56\)-page PDF and the canonical Overleaf workbench both compile with
  zero errors, warnings, undefined references, or undefined citations.
\item Synchronized the verified analytic Picard proof, literature references,
  status files, dependency/morphism/uncertainty ledgers, and research log to
  the canonical validate/repair/disprove project.  The current workbench
  distinguishes \Verified\ (a complete auditable proof is written here) from
  \Open\ (not yet settled); neither label is a claim of community acceptance.
\item Rebuilt and synchronized the public annotated edition.  Its current
  \(123\)-page PDF places the independent review record first, the subsequent
  Codex results (including Proposition~10.8) in the front matter, the
  \(108\)-page reconstruction next, and the indexed historical appendix last.
  Both canonical public compile checks report zero errors and zero warnings.
\item The optional cellular-collapse/Mayer--Vietoris comparison and the
  external CDP adjudication remain \Open.  Historical reviewer wording is
  retained and labelled rather than silently rewritten.
\end{itemize}

\subsection*{Round 10 --- 26 August 2026: branch reconciliation and single-project consolidation}

\begin{itemize}[leftmargin=1.8em]
\item Reconciled the later 18-page analytic branch with the full 56-page
workbench.  Its conservative topology labels were a fork checkpoint and do
not supersede the complete intrinsic nearby-cycle/Leray and based van Kampen
proofs already present here.  In particular, V1, the three Leray
transgressions, $H^2(X;\mathbb Z)=0$, and sphere recognition remain
\Verified; the cellular-collapse comparison alone remains \Open.
\item Imported the branch's genuinely independent improvement: the cusp
toric volume form proves $\omega_X$ trivial near $W$, and every
$L\in\operatorname{Pic}^0(X)$ pulls back trivially to the
$\mathrm{dP}_6$ normalization.  The conductor, compact-Cartier
Ramis--Ruget duality, and top-degree base-change chain now proves
$R^2f_*(T_X\otimes L)\ne0$ on $\operatorname{Pic}^0(X)$ without any global
topology.
\item Isolated the exact all-line-bundle obstruction
$r=\eta^*i^*:\operatorname{Pic}(X)\to\operatorname{Pic}(\mathrm{dP}_6)$.
The uniform trivial-pullback hypothesis is equivalent to $r=0$.  For the
constructed candidate this is already verified from the independently proved
$H^2(X;\mathbb Z)=0$; only a second, transition-function proof of $r=0$
without topology remains open.
\item Added the direct Ramis--Ruget proof of degree-two duality for a compact
Cartier fibre in a possibly noncompact smooth ambient threefold, together
with explicit dualizing-complex, adjunction, trace, compact-support,
finiteness, and variance checks.
\item Recorded the all-degree terminal homology implication as a complete
finite-algebra proposition while retaining the explicit warning that the
printed terminal complexes must still be identified with geometric maps.
The canonical workbench separately supplies that identification through its
verified intrinsic route.
\item Preserved the canonical closed equalizer proof of Proposition~4.7(iv)
and rejected the shorter branch's stale reopening of it.  Claim identifiers
were merged semantically rather than appended blindly, avoiding collisions
between abstract CDP theorems and candidate-specific Pic$^0$ applications.
\item Moved active publication work into the single canonical 123-page
project.  The former workbench project remains intact as a compiled
checkpoint with a migration notice; no sharing controls were changed.
\end{itemize}

\subsection*{Round 11 --- 26 August 2026: Reviewer B analytic follow-up}

\begin{itemize}[leftmargin=1.8em]
\item Audited Reviewer~B's seven-page candidate-specific validation update
against the current workbench.  Its genuinely new exposition is the explicit
finite-complex universal-coefficient calculation: for every vector bundle in
a proper flat analytic family of scheme-theoretic surfaces,
\[
R^qf_*E=0\ (q>2),\qquad
(R^2f_*E)_b\otimes\mathbb C(b)\xrightarrow{\sim}H^2(X_b,E_b)
\]
at every fibre, including a nonreduced fibre.  This proof is now incorporated
as \cref{cor:surface-top-base-change} and bound to the precise analytic
source already cited in the workbench.
\item The follow-up independently confirms, from the named inputs CAND-X,
CAND-W, and TOP-H2, the normalization-trivial twist, conductor-supported
ambient differential, compact-Cartier duality, and
$R^2f_*(T_X\otimes L)\ne0$ for every $L\in\operatorname{Pic}(X)$.  Its
OPEN labels for those three inputs record its own cutoff and do not downgrade
the later construction and topology proofs in this workbench.
\item Reran its eight finite checks successfully.  They certify $B_0$, the
$xy/xyz$ conormal annihilators, finite conductor inclusion--exclusion
templates, conductor-ideal membership of $d(xy)$ and $d(xyz)$, and $p=-1$;
they do not substitute for the global candidate inputs.
\item Corrected the typed dependency ledger: Proposition~10.8 is not an
antecedent of the all-Pic direct-image nonvanishing.  That conclusion uses the
cusp conductor, duality, top-degree base change, and the restriction map
$r=\eta^*i^*$ killed by the independently proved $H^2(X;\mathbb Z)=0$.
\item Preserved the pristine 23-file Reviewer~B archive and an Internet-safe
curated evidence selection.  The report's shorthand ``SNC'' is not adopted
as current terminology: $W$ is irreducible and nonnormal, with the stated
normal-crossings local models, not a simple-normal-crossings divisor.
\item The single-project local tree now contains the full workbench.  Earlier
log references to remote synchronization describe prior checkpoints; the
current consolidated Overleaf tree is not recorded as synchronized until its
root file tree and remote compile have been inspected directly.
\end{itemize}

\subsection*{Round 12 --- 27 August 2026: collapse and specialization closed}

\begin{itemize}[leftmargin=1.8em]
\item Re-audited the complete candidate-specific collapse chain in the
reconstructed packets, rather than promoting the manuscript's conclusion by
reference.  The straightening homeomorphism is continuous at the toric
boundary; the polar quotient has exactly the printed connected stabilizers;
the positive locus and its honeycomb boundary are contractible manifolds; and
the compact quotient, collar, common deck group, lifted deformation, and
phase-spreading arguments give a genuine equivariant strong deformation
retraction of the cusp tube onto the singular fibre.  Obligation C1 is now
\Verified.
\item Modified the end map equivariantly so that its restriction to a marked
smooth fibre is the explicit honeycomb collapse.  Direct cellular and
intersection calculations give, in every degree, the exact integral sequence
\[
0\longrightarrow\operatorname{im}(\textstyle\bigwedge^qM_0-I)
\longrightarrow\textstyle\bigwedge^q\Lambda
\xrightarrow{c_*}H_q(W;\mathbb Z)\longrightarrow0.
\]
The degree-two and degree-three matrices, their primitive kernels, saturation,
and torsion-freeness are written out in
\cref{prop:verified-collapse-homology}.  Obligation C2 is now \Verified.
\item Compared the topological collapse with the analytic nearby-cycle unit,
not merely with its rank.  For a sufficiently small tube and radial marking,
\[
\operatorname{sp}^q
=i_F^*(i_W^*)^{-1}=i_F^*r^*=c^*
:H^q(W;\mathbb Z)\longrightarrow H^q(F;\mathbb Z)
\]
for every degree.  This is an exact integral equality, independent of the
chosen lift because the image is monodromy-invariant.  Obligation C3 is now
\Verified.
\item Checked the topological inputs against Brown's collar theorem, the
nearby-cycle realization against Dimca's proper-specialization formalism, and
Clemens's degeneration paper as historical comparison.  These sources supply
general tools only; the candidate-specific retraction and comparison are
proved in the workbench.
\item Kept MV1 and the Mayer--Vietoris comparison part of L1 \Open.  The
remaining task is to orient the actual global collars, derive the inclusion
maps $(i_{U*},-i_{N*})$ on primitive cycles, and exhibit a commuting integral
comparison with the verified Leray transgressions.  Equality of terminal
Smith forms does not discharge that task.
\end{itemize}

\subsection*{Round 13 --- 27 August 2026: global Mayer--Vietoris and Leray comparison closed}

\begin{itemize}[leftmargin=1.8em]
\item Corrected the earlier audit's one-shot-cover conflation.  The printed
$\alpha_q$ belong to the finite-band decomposition
$X^\circ=X_1\cup X_2$ and are
$\alpha_q=(\pi_{1*},-\pi_{2*})$; the cusp is attached in the second
decomposition $X=N'_0\cup X^\circ$ through
$K_q=\ker(H_q(M)\to H_q(W))$.
\item Used the verified finite radial markings to identify the actual band
inclusions with $\pi_j$, then computed $\alpha_0,\ldots,\alpha_4$ in named
primitive bases.  The degree-one and degree-two matrices have Smith form
$\operatorname{diag}(1,1,1,0)$; the degree-two left kernel is the primitive
functional $(4,2,3,2)$; and the top map is $(3,-4)$ with the complex
covering orientations.
\item Derived the cusp kernels from literal cycles:
\[
K_1=\mathbb Z[\widetilde a_0],\qquad
K_2=\mathbb Z\langle\xi_{\widehat w},
                         \xi_{\widehat\delta}\rangle.
\]
The actual global degree-one block has determinant $-p$, and the reduced
degree-two block is $\operatorname{diag}(-p,1)$.  A separate one-shot
$9\times9$ collar matrix has determinant $-p$ for the candidate family.
At $p=-1$ all these blocks are unimodular.
\item Built the missing comparison diagrams.  On the clockwise
disc--pair-of-pants cover, the local discrepancy tuple is the negative of the
transpose-collar tuple, so the degree-one square commutes.  Sweeping by the
global $S^1_{\widehat\delta}$-action gives the degree-two square, and the
Leray derivation rule fixes degree three.  In the resulting primitive bases,
\[
d_2(12\gamma)=-p\omega,\quad
d_2(2q)=-p\omega[\delta],\quad
d_2(2\gamma uw)=-p\omega[\gamma\delta].
\]
Thus MV1 and L1 are now \Verified, not merely numerically consistent.
\item Two independent read-only audits reconstructed the same maps and
attempted to falsify the primitive bases, signs, and indices.  Exact integer
arithmetic independently confirmed the two Smith forms, the primitive left
kernel, $\det D_1=-p$, and the one-shot determinant
$1-12\ell_0=-p$ for $(\ell_1,\ell_2)=(1,-1)$.
\item The annotated edition and its historical verdicts remain frozen as
provenance.  This new mathematics is recorded only in the live workbench and
its research ledgers.
\end{itemize}

\subsection*{Round 14 --- 27 August 2026: exact Hodge boundary and a failed shortcut}

\begin{itemize}[leftmargin=1.8em]
\item Reconstructed the candidate-specific holomorphic-form vanishing and
combined it with the proved structure-sheaf direct images, Serre duality,
the sphere topology, and Hirzebruch--Riemann--Roch.  With
$m=h^{1,2}(X)$, the only undetermined Hodge entries are
\[
h^{1,2}=h^{2,1}=m,\qquad h^{1,1}=h^{2,2}=m+1.
\]
\item Determined the entire Fr\"olicher sequence.  Its two middle $d_1$ rows
are exact with ranks $(1,m)$ and $(m,1)$, and no higher differential can
survive; hence $E_2=E_\infty$ and $\dim E_2=2$ for every $m\geq0$.
\item Computed the period-derivative determinant
\[
K=\det(\dot Z)=6\dot\mu^{\,2}-\dot\tau\dot\beta.
\]
The toric cusp normal form gives $K\neq0$ on a punctured cusp neighbourhood,
so the Kodaira--Spencer map has generic rank two.  At the order-three
elliptic point, however,
$\dot\mu=-\dot\tau/3$ and $\dot\beta=2\dot\tau/3$, so $K=0$.  This refutes
the proposed passage from a nowhere-zero matrix to an everywhere-invertible
one.
\item Reduced the unknown exactly to
\[
m=h^1(\mathbb P^1,R^1f_*\Omega_X^1)
  +h^0(\mathbb P^1,R^2f_*\Omega_X^1).
\]
The generic ranks are $(1,3,2)$ in direct-image degrees $0,1,2$.  Determining
$m$ requires the splitting degrees and torsion created by the two cyclic
multiple fibres and the toric cusp.  The expected value $m=1$ remains
\Open.
\item Checked the scope against Ugarte and Angella.  Their general Hodge
constraints agree with the resulting table but do not supply the missing
candidate-specific singular-fibre direct images.  The public Angella source
was used for audit and is cited rather than redistributed.
\item In accordance with the current provenance directive, no annotation,
overview, or reconstructed-manuscript file was changed; all new mathematics
was confined to the live workbench and its research ledgers.
\end{itemize}

\subsection*{Round 15 --- 27 August 2026: the Hodge calculation closed}

\begin{itemize}[leftmargin=1.8em]
\item Replaced the generic-rank reduction by the actual logarithmic relative
cotangent calculation.  For
\(
\mathcal E=\Omega_X^1(\log D)/f^*\Omega_B^1(\log\Sigma)
\), the two cyclic invariant lattices and the toric cusp frame give
\[
f_*\mathcal E=\mathcal O\oplus\mathcal O(-1),\quad
R^1f_*\mathcal E=\mathcal O^2\oplus\mathcal O(-1)^2,\quad
R^2f_*\mathcal E=\mathcal O\oplus\mathcal O(-1).
\]
The variable tangent and cotangent generators pair to $t$, not to a unit;
this resolves the apparent contradiction with
$\mathcal E^\vee\simeq T_{X/B}$ and refutes the shortcut that direct image
commutes with duality across a multiple fibre.
\item Computed the logarithmic Kodaira--Spencer determinant as a section of
$\mathcal O(2)$.  Its invariant exponents force one zero at each finite
point, while it is a unit at the cusp, hence
\[
\operatorname{div}(\det\kappa_0)=p_1+p_2.
\]
The local exterior matrix
$(-C,A,-D,B)$ proves that $\kappa_1$ is sheaf-surjective.
\item Matched the logarithmic residue sequence through all three fibres.
The two skyscraper terms $\operatorname{coker}\kappa_0$ map isomorphically
to $H^1(S_1,\mathcal O)\oplus H^1(S_2,\mathcal O)$, while the weighted
degree-zero residue has kernel $\mathcal O(-2)$.  This yields the absolute
cotangent direct images up to the single global rank of $H^0(\kappa_1)$.
\item Closed that scalar without an Abel--Jacobi assumption.  At the
order-four lift, $\tau-i$ has exact order two, so $\dot\tau(z_2)=0$; the
simple logarithmic determinant zero forces raw $K(z_2)\ne0$, and therefore
$\dot\mu(z_2)\ne0$.  The normalized global form then has nonzero image under
$H^0(\kappa_1)$.  Consequently
\[
R^0f_*\Omega_X^1=\mathcal O(-2),\quad
R^1f_*\Omega_X^1=\mathcal O\oplus\mathcal O(-1)^2,\quad
R^2f_*\Omega_X^1=\mathcal O\oplus\mathcal O(-1),
\]
and
\[
h^{1,2}=1,\qquad h^{1,1}=2.
\]
Three independent audits converged on the scalar after one proposed
$h^{1,2}=2$ shortcut was falsified: it had divided by $\dot\tau$ at the
order-four point, where $\dot\tau$ vanishes and the quotient has a pole.
\item Corrected the manuscript's local wording in the live workbench:
$\Omega^1_{X/B}$ is torsion-free but non-locally-free at the cusp's
$xy/xyz$ strata; its torsion occurs at the two multiple fibres.  The frozen
source statement remains unchanged as provenance.
\end{itemize}

\subsection*{Round 16 --- 27 August 2026: Bott--Chern and Aeppli diamonds}

\begin{itemize}[leftmargin=1.8em]
\item Audited McHugh's complete compact-threefold calculation in the primary
published record and the public version-four LaTeX source
\cite[Theorem~6.3, Proposition~6.4, equation~(7), \S7.1]{McHugh2019}.
Its load-bearing
identity is
\[
 h^{1,2}_{BC}=h^{1,2}_{\bar\partial}
 +\dim_{\mathbb C}E_2^{3,1}-k^{2,0},
 \qquad k^{p,q}=\dim K^{p,q}\le b_{p+q}.
\]
For this candidate the three terms on the right are $1,0,0$.
\item Substituted the complete Hodge diamond, $E_2=E_\infty$, and sphere
cohomology into McHugh's threefold table.  This gives the Bott--Chern rows
$q=3,2,1,0$
\[
 (0,0,1,1),\quad(0,1,2,1),\quad(0,2,1,0),\quad(1,0,0,0),
\]
and, by Schweitzer's Hodge-star duality
\cite[\S2.c, pp.~9--10]{Schweitzer2007}, the Aeppli rows
\[
 (0,0,0,1),\quad(0,1,2,0),\quad(1,2,1,0),\quad(1,1,0,0).
\]
The total-degree non-$\partial\bar\partial$ vector is
$(0,2,4,4,4,2,0)$.
\item At this Round~16 checkpoint every numerical Bott--Chern and Aeppli
dimension was \Verified, while explicit period-coordinate representatives
and the product structure were still \Open\ strengthenings.  Round~23 below
records their subsequent verification; they were never premises of the
Round~16 dimension calculation.
\item A fresh adversarial audit independently replayed the finite invariant
exponents, elementary modifications, logarithmic Kodaira--Spencer divisor,
residue cancellation, order-four scalar, coherent Leray calculation, and
the McHugh specialization.  It found no fatal gap and identified several
compressed explanations; those explanations were expanded in the proof.
\end{itemize}

\subsection*{Round 17 --- 27 August 2026: adversarial quotient and cusp-map audit}

\begin{itemize}[leftmargin=1.8em]
\item Corrected the mechanism behind $k^{p,q}\le b_{p+q}$.  There is in
general no well-defined injection $K^{p,q}\hookrightarrow
H^{p+q}_{\mathrm{dR}}$.  If $Z^{p,q}=\ker\partial\cap\ker\bar\partial$ and
$B_K^{p,q}$ is McHugh's denominator, then
\[
 \ker\bigl(Z^{p,q}\to H^{p+q}_{\mathrm{dR}}\bigr)\subset B_K^{p,q}.
\]
Therefore the de Rham image of $Z^{p,q}$ surjects onto $K^{p,q}$, which is
the correct proof of the inequality.  The sphere-cohomology specialization,
$k^{2,0}=0$, and all Bott--Chern/Aeppli numbers are unchanged.
\item Expanded the cusp cotangent certificate.  The normalization quotient
$Q=\nu_*\mathcal O_{\mathrm{dP}_6}/\mathcal O_W$ has
$h^\bullet(Q)=(2,1,0)$; the branch-difference map sends
\[
 d\log x_1\mapsto(-1,0,1),\qquad
 d\log x_2\mapsto(0,-1,-1),
\]
and is an isomorphism onto $H^0(Q)$.  In bases transported by the connecting
isomorphisms, the next map is the explicit surjective row
$(0,1,-1,0):\mathbb C^4\to H^1(Q)$.  Hence the three cusp direct-image
modules are $t\mathcal O^2$, the kernel of that row modulo $t$, and
$\mathcal O^2$, respectively.  This makes the earlier rank-only remark
auditable and confirms that the cusp is a non-locally-free, not a torsion,
locus of $\Omega^1_{X/B}$.
The sheaf exactness is also checked before taking cohomology: at $t=xy$ the
transverse cokernel is $R/(x,y)$, while at $t=xyz$ the ordinary image in
$Re_1\oplus Re_2$ is generated by $xe_1,ye_2,z(e_1+e_2)$, exactly the local
normalization-quotient presentation.
\item Closed the primary-source layer for the general analytic tools used by
the calculation: Cartan for finite analytic quotients, Deligne for logarithmic
conormal/residue sequences, Katz--Oda for connecting maps and exterior
products, Grothendieck and the Stacks Project for modifications and splitting
on $\mathbb P^1$, and Fr\"olicher for the spectral-sequence identifications.
These citations support the formalism; the invariant exponents, cusp rows,
period determinant, and order-four scalar remain proved directly for this
candidate.
\end{itemize}

\subsection*{Round 18 --- 28 August 2026: Leray-edge all-Pic restriction}

\begin{itemize}[leftmargin=1.8em]
\item Separated two logically different questions that had been compressed
into the sentence ``$H^2(X;\mathbb Z)=0$ kills the restriction.''  For every
$L\in\operatorname{Pic}(X)$, the degree-two Leray edge section satisfies
\[
 e(c_1(L))\in E_\infty^{0,2}\subseteq\ker d_2^{0,2}=0,
\]
because the already verified candidate calculation gives
$d_2^{0,2}(2q)=\omega[\delta]$.  Functoriality of the filtered spectral
sequence, analytic proper base change, and naturality of the exponential
boundary identify evaluation at $p_0$ with $c_1(L|_W)$.  Thus
$c_1(\eta^*i^*L)=0$, and injectivity of $c_1$ on
$\operatorname{Pic}(\mathrm{dP}_6)$ proves
$\eta^*i^*L\simeq\mathcal O$.
\item Audited the dependency order.  The proof cites
\cref{lem:higher-transgressions}, which is proved before and without
\cref{thm:h2-vanishing}; there is no logical circularity.  It avoids invoking
the packaged $H^2(X;\mathbb Z)=0$ theorem, but it is not independent evidence:
the same $d_2^{0,2}$ is the decisive input in that theorem.
\item Checked the formal literature layer in the local Stacks LaTeX corpus:
Tags~01F2, 011N, 011O, 012O, 012U, and~012W give Leray, permanent-cycle,
functoriality, associated-graded, and convergence statements.  Dimca's
Theorem~2.3.26 supplies analytic proper base change.  The evaluation/Chern
class compatibility is written as an explicit commutative square in
\cref{prop:leray-edge-restriction}, not left as an unstated convention.
\item Kept the stronger transition-function question open.  A proof of
$\eta^*i^*=0$ from candidate transition functions or the relative Picard
functor, independent of the integral Leray transgressions, would be a second
certificate.  The verified all-Pic conductor and $R^2$ nonvanishing results
do not depend on obtaining that second certificate.
\end{itemize}

\subsection*{Round 19 --- 28 August 2026: direct boundary-Picard certificate}

\begin{itemize}[leftmargin=1.8em]
\item Replaced the open ``second certificate'' from Round~18 by a complete
direct boundary calculation.  For $L\in\operatorname{Pic}(X)$, put
$E=c_1(L|_F)$.  On the actual pair-of-pants core, subtracting the invariant
vertical form gives a crossed homomorphism with coinvariant defects
$\bar d_1,\bar d_2,\bar d_0$ satisfying
\[
 \bar d_1+\bar d_2+\bar d_0=0.
\]
The order-four specialization lattice forces $E=2nq$.  The order-three and
order-four norm identities, with the proved logarithmic conjugacies, give
\[
 \bar d_1=4n[\delta],\qquad
 \bar d_2=-3n[\delta].
\]
The zero-section $\widehat\delta$-sweep at the cusp is the boundary of the
explicit solid torus
\[
 \{|u|=\rho,\ x_1=1,\ |x_2|\le1\},
\]
so $\bar d_0=0$.  Since $V_G=\mathbb Z[\delta]$, balance gives $n=0$.
Naturality and injectivity of local integral specialization then yield
$c_1(L|_W)=0$, and injectivity of $c_1$ on
$\operatorname{Pic}(\mathrm{dP}_6)$ gives
$\eta^*i^*L\simeq\mathcal O$.
\item Made every bridge explicit: the crossed-homomorphism representative
independence, the fibre-first/base-first sign, the finite-cover norm and the
direction of the logarithmic conjugacy, the order-four parity step, the cusp
boundary parametrization, and the nearby-cycle Chern-class naturality
square.  Two independent read-only proof audits returned \textsc{promote}.
\item Added the exact arithmetic harness
\Source{supporting_materials/reproducibility/verify_direct_boundary_picard.py}.
It passes all fourteen lattice, contraction, parity, and primitive-obstruction
checks; its scope note correctly leaves the geometric bridge lemmas to the
written proof.
\item Promoted this direct boundary route to the primary dependency of the
all-Pic conductor/direct-image theorem.  The Leray-edge proof remains an
independent certificate.  A still stronger proof using only
$\mathcal O^*$ cocycles, support functions, or the relative Picard functor,
with no local integral specialization or bounding chain, remains \Open.
\item Corrected the smooth-recognition literature chain: Smale's 1961
Theorem~A supplies homeomorphism, while Smale's 1962 Theorem~1.1 and
Corollary~1.2 supply the direct diffeomorphism statement in dimension six;
equivalently, the Kervaire--Milnor table and following remark give
$|\Theta_6|=1$ and identify its elements with diffeomorphism classes.
\end{itemize}

\subsection*{Round 20 --- 28 August 2026: independent Hodge-map audit}

\begin{itemize}[leftmargin=1.8em]
\item Two independent read-only audits reconstructed the cyclic characters
at both finite fibres, the invariant bases (11.4), the evaluation elementary
divisors, the residue cancellation, the cusp branch-difference maps, and the
order-four scalar.  Both found the conclusions
\[
 R^1f_*\Omega_X^1\simeq\mathcal O\oplus\mathcal O(-1)^{\oplus2},
 \qquad h^{1,2}=1,\qquad h^{1,1}=2
\]
unchanged.
\item Repaired three exact-wording defects.  The determinant formula (11.9)
now includes its nowhere-zero frame unit; the constant cusp
Kodaira--Spencer matrices are stated modulo $t$ rather than as full germs;
and the global fixed cotangent line is introduced only after constructing
and splitting the evaluation sequence
\[
 0\to\mathcal O(-1)\to f_*\mathcal E
 \xrightarrow{\operatorname{ev}_\xi}\mathcal O\to0.
\]
The order-four value is likewise stated in a normalized determinant frame,
with arbitrary frames differing by a nonzero scalar.
\item Split the former cusp morphism record into the two actual maps
$\rho_0,\rho_1$ and the following long-exact-sequence connecting maps.
The local $xy/xyz$ injections also receive their own \Refuted\ ledger
entry against the manuscript's claim that torsion occurs at all three
singular fibres.  Near the cusp $\Omega^1_{X/B}$ is torsion-free but
non-locally-free; after pullback to $t=s^{m_j}$ at a finite fibre its torsion
is generated by $ds$ and killed by $s^{m_j-1}$.
\item Rechecked the Bott--Chern substitution against McHugh's public
version-four LaTeX source and the logarithmic formalism against the primary
Cartan, Deligne, and Katz--Oda records.  The numerical formulas and
candidate-specific use are unchanged.  The literature ledger now identifies
Hartshorne's cited book correctly as \emph{Algebraic Geometry}, Chapter~V,
and labels Reviewer~B separately as audit provenance rather than literature.
\end{itemize}

\subsection*{Round 21 --- 28 August 2026: status and morphism-ledger normalization}

\begin{itemize}[leftmargin=1.8em]
\item Normalized the methodology to the four literal statuses used by every
live surface: \textsc{verified}, \textsc{conditional}, \textsc{open}, and
\textsc{refuted}.  Repair is recorded as an operation: a false printed claim
is \textsc{refuted}, while its separately proved replacement is
\textsc{verified}.  In particular, \textsc{conditional} now appears only for
a proved implication that names the established or assumed hypothesis on
which its conclusion depends.
\item Removed the last sign ambiguity from the first-transgression morphism
record.  With the common orientation used in the comparison theorem, the
general coefficient is
\[
 d_2(12\gamma)=-p\,\omega,
\]
and hence equals \(\omega\) for the constructed candidate \(p=-1\).
\item Retyped the Bott--Chern quotient morphism.  From
\[
 \ker\bigl(Z^{p,q}\longrightarrow H^{p+q}_{\mathrm{dR}}(X;\mathbb C)\bigr)
 \subseteq B_K^{p,q}
\]
one obtains the well-defined surjection
\[
 \operatorname{im}\bigl(Z^{p,q}\to H^{p+q}_{\mathrm{dR}}(X;\mathbb C)\bigr)
 \longtwoheadrightarrow K^{p,q},
 \qquad [\alpha]_{\mathrm{dR}}\longmapsto[\alpha]_K.
\]
This is the precise map yielding \(k^{p,q}\le b_{p+q}\); no injection of
\(K^{p,q}\) into de Rham cohomology is asserted.
\item Recorded the two Kodaira--Spencer connecting maps with their separate
types: \(\kappa_0:f_*\mathcal E\to\mathcal L\otimes R^1f_*\mathcal O_X\)
and
\(\kappa_1:R^1f_*\mathcal E\to\mathcal L\otimes R^2f_*\mathcal O_X\).
The ledger no longer certifies the second map under the first map's domain.
\item Separated three further bundled records into the geometric maps that
are actually proved.  The finite-band Mayer--Vietoris map
\(\alpha_q:H_q(F)\to H_q(S_1)\oplus H_q(S_2)\) is now distinct from the
cusp-attachment map
\(\Delta_q=(c_*,-j_*):H_q(M)\to H_q(W)\oplus H_q(X^\circ)\); the four
finite-filling and cusp inclusions used by van Kampen are recorded
individually; and the direct boundary-Picard certificate is typed as
\(r=\eta^*i^*:\operatorname{Pic}(X)\to\operatorname{Pic}(\mathrm{dP}_6)\).
The degree-two cusp block is written with its proved orientation,
\(j_*\xi_{\widehat\delta}=-p\,e\), rather than an ambiguous punctuation
artifact; the same correction and the missing \(\backslash\texttt{longmapsto}\)
command were propagated through the comparison square (L.7).
\end{itemize}

\subsection*{Round 22 --- 29 August 2026: theorem-status semantics}

\begin{itemize}[leftmargin=1.8em]
\item Reclassified three complete implication theorems from
\textsc{conditional} to \textsc{verified}: smooth six-sphere recognition,
the abstract topology application, and the conductor/direct-image theorem.
Their displayed hypotheses are theorem antecedents whose implications are
proved; none imports an unproved external hypothesis.  The live claim counts
are therefore \(38\) \textsc{verified}, \(0\) \textsc{conditional},
\(5\) \textsc{open}, and \(4\) \textsc{refuted}.
\item Retained \textsc{conditional} as an available status with its literal
meaning: it is reserved for a conclusion whose proof imports a named
unproved or assumed hypothesis, such as a theorem conditional on a
conjecture.  No current claim in the workbench has that status.
\item Replaced the stale description of the CDP monodromy step as an
\emph{unresolved dependency} by its proved disposition: the direct-product
homology step is \textsc{refuted} by the semidirect-product calculation and
must be replaced by coinvariants.  External replication remains welcome but
is not a hypothesis or a gate on continuing the proof.
\end{itemize}

\subsection*{Round 23 --- 29 August 2026: local Picard boundary and Bott--Chern algebra}

\begin{itemize}[leftmargin=1.8em]
\item Promoted the explicit local cusp calculation to a verified claim.  The
equivariant Cartier divisor
\[
 \mathcal D_Q=\sum_{v\in\mathbb Z^2}(v_1^2-6v_2^2)E_v
\]
descends to $\mathcal L_Q\in\operatorname{Pic}(N_0)$ with generic class
$-2q$ and normalized central degree vector
$(-12,2,0,-12,2,0)$.  Thus
$\operatorname{Pic}(N_0)\to\operatorname{Pic}(\mathrm{dP}_6)$ is nonzero.
This does not conflict with the two verified zero maps whose domain is
$\operatorname{Pic}(X)$; it constrains, but does not close, the optional
global purely analytic route.
\item Closed the explicit-representative and product obligation.  Global
$\bar\partial$-closed representatives $G,Q$ have exact marked-fibre
period-coordinate classes and yield the nine positive-degree Bott--Chern
generators
$\mathsf A,\mathsf B,\mathsf C,\overline{\mathsf C},
\mathsf D,\overline{\mathsf D},
\mathsf E,\overline{\mathsf E},\mathsf T$.
For their span $N$ the complete algebra is
\[
 H^{*,*}_{BC}(X)=\mathbb C\mathbf1\oplus N,
 \qquad N^2=0.
\]
The proof supplies classes, exact fibre formulas, and every product, but not
a preferred literal global coefficient atlas for the noncanonical smooth
forms through the two finite fillings and cusp.
\item Preserved the stable closure ID
\texttt{S6-OPEN-HODGE-REPRESENTATIVES} and changed its status in place from
\textsc{open} to \textsc{verified}.  After the local Picard addition and this
closure, the live claim ledger has \(40\) \textsc{verified}, \(0\)
\textsc{conditional}, \(4\) \textsc{open}, and \(4\) \textsc{refuted}
claims.  The uncertainty ledger has \(10\) \textsc{verified}, \(5\)
\textsc{open}, and \(3\) \textsc{refuted} entries; the dependency ledger has
\(92\) rows; and the morphism ledger has \(25\) verified typed maps.
\end{itemize}

\subsection*{Round 24 --- 30 August 2026: corrected CDP theorem and cusp normalization}

\begin{itemize}[leftmargin=1.8em]
\item Acquired and read the primary arXiv \LaTeX\ sources for CDP20
(arXiv:1904.11179) and CDP98 (arXiv:math/9607215), in addition to the
published PDFs already audited.  The source confirms verbatim that CDP20
Proposition~2.4 is a sufficient relative-vanishing criterion logically
separate from the two failed attempts to establish its premise.
\item Proved the strongest intact replacement theorem.  If one
\(L_0\in\operatorname{Pic}(X)\) has
\[
 H^2(X,T_X\otimes L_0)=0,
\]
then generic endpoint vanishing and semicontinuity give
\(H^i(X,T_X\otimes M)=0\) for \(i\ne1\), and HRR gives
\[
 \chi(X,T_X\otimes M)=\tfrac12c_3(X)\le0.
\]
The relative condition \(R^2f_*(T_X\otimes L)=0\) for one \(L\) implies
this hypothesis by projection formula, Serre vanishing on the base, and the
two-column Leray spectral sequence.
\item Replaced the invalid p.~692 reduction by an exact reduced-fibre
theorem.  After the two endpoint vanishings, ambient cotangent vanishing is
equivalent to injectivity of
\[
 \delta_c:H^0(W_c,\operatorname{tors}\Omega^1_{W_c}\otimes A_c)
 \longrightarrow H^1(W_c,N^*_{W_c/X}\otimes A_c).
\]
The criterion is imposed on the whole reduced Cartier fibre, including a
reducible one; nonreduced fibres require separate nilpotent-thickening data.
Cartier duality, top-degree base change, and Nakayama then prove the relative
vanishing.  The candidate fails both repaired hypotheses by distinct proofs:
\(c_3(X)=2\) and the corrected global theorem give
\(H^2(X,T_X\otimes L_0)\ne0\) for every \(L_0\) by contraposition, while the
conductor calculation gives \(R^2f_*(T_X\otimes L)\ne0\) for every \(L\).
Thus this repair does not restore the refuted original theorem.
\item Made the monodromy repair exact: the CDP98 restriction diagram reaches
only the invariant integral lattice, and its last linear step requires that
this lattice span \(H^1(F,\mathcal O_F)\).  Trivial integral monodromy is a
stronger sufficient hypothesis; neither condition follows from the original
CDP20 assumptions.
\item Expanded the modular-cusp contradiction in
\cref{prop:period-input-audit}.  A constant modular conjugation first carries
the relevant inverse-\(j\) component's boundary cusp to infinity; the
extended \(q\)-coordinate cannot have modulus one by the maximum principle,
cannot be nonzero because \(j\) has a pole, and cannot vanish while possessing
a single-valued logarithm.  This closes the compressed normalization step
identified by the adversarial upstream-geometry audit.
\item Preserved the stable ID \texttt{S6-OPEN-CDP-REPAIR} and changed its
status from \textsc{open} to \textsc{verified}.  The live ledgers now contain
\(41\) \textsc{verified}, \(0\) \textsc{conditional}, \(3\) \textsc{open},
and \(4\) \textsc{refuted} claims; \(11\) \textsc{verified},
\(4\) \textsc{open}, and \(3\) \textsc{refuted} uncertainties; \(96\)
dependency rows; and \(25\) verified typed morphisms.
\end{itemize}

\subsection*{Round 25 --- 30 August 2026: purely analytic all-Picard certificate}

\begin{itemize}[leftmargin=1.8em]
\item Closed the optional third proof of
\(r=\eta^*i^*:\operatorname{Pic}(X)\to\operatorname{Pic}(\mathrm{dP}_6)\).
For $E=c_1(L|_F)=kq$, logarithms of the actual $\mathcal O^*$ factors define
a crossed defect whose finite norms and cusp winding satisfy
\[
 (\bar d_1,\bar d_2,\bar d_0)
   =(2k,-\tfrac32k,0)[\delta],
 \qquad
 \bar d_1+\bar d_2+\bar d_0=0.
\]
The cusp value is computed by the argument principle in an extending Stein
chart; it is not a Chern evaluation on a bounding chain.  Hence $k=0$.
\item Proved the remaining central-fibre implication analytically.  For the
three independent normalized central degrees,
\[
 B=\begin{pmatrix}d_2+d_3&d_3\\d_3&d_1+d_3\end{pmatrix},
 \qquad
 E_L((\lambda,0),(0,m))=-B(B_0\lambda,m).
\]
Since $B_0$ is unimodular, $E_L=0$ forces $B=0$, all six central degrees
vanish, and $r(L)$ is trivial.  This proof uses neither global integral
Leray nor local integral nearby-cycle specialization.
\item Repaired three defects found during adversarial audit: the logarithmic
and linear finite-filling lifts are now distinguished and transported by the
actual punctured-filling conjugacy; the cusp torus uses
$b=a_u$, $f=m_{x_2}-a_u$ rather than a false fixed-base bounding assertion;
and frame invariance is proved in the underlying commuting deck group rather
than by assuming commutation of central-extension lifts.  The final
adversarial disposition was a pass.
\item The exact verifier passed all $25/25$ checks, including the six support
values, six central degrees, polarization matrix, mixed block, explicit
$Q=x^2-6y^2$ example, $\det B_0=1$, and
$2-\tfrac32=\tfrac12$.  The live ledgers now contain $42$ \textsc{verified},
$0$ \textsc{conditional}, $2$ \textsc{open}, and $4$ \textsc{refuted}
claims; $12$ \textsc{verified}, $3$ \textsc{open}, and $3$ \textsc{refuted}
uncertainties; $100$ dependency rows; and $26$ verified typed morphisms.
\end{itemize}

\subsection*{Round 26 --- 30 August 2026: exact-normalization supersession}

\begin{itemize}[leftmargin=1.8em]
\item Fixed the finite-generator signs as one simultaneous triangle-group
orientation choice.  With \(g_0=(g_1g_2)^{-1}\),
\[
 g_1'(z_1)=e^{-2\pi i/3},\qquad g_2'(z_2)=-i,\qquad
 g_0\text{ parabolic}.
\]
For
\[
 S_1(w)=\frac{w-1}{w},\qquad S(w)=-\frac1w,\qquad P(w)=w-1,
\]
the exact modular deck certificate is
\[
 \varphi(g_1)=S_1,\qquad\varphi(g_2)=S,\qquad
 \varphi(g_0)=P,\qquad S_1SP=1.
\]
The quadratic ramification at \(z_2\) yields only
\((g_2'(z_2))^2=-1\); retaining the chosen \(g_1\) orientation and the
primitive parabolic cusp class forces the sign \(-i\).

\item Made the invariant fibre class and its orientation convention literal.
On a marked fibre put
\[
 y=\operatorname{Im}\tau,\qquad m=\operatorname{Im}\mu,\qquad
 r=\operatorname{Im}\beta,\qquad
 \Delta=r-\frac{6m^2}{y}<0,
\]
\[
 \Gamma=\frac{i}{2\Delta}
 \left(d\bar\zeta_2-\frac{m}{y}d\bar\zeta_1\right),\qquad
 U=\frac{i}{2y\Delta}
 \left(r\,d\bar\zeta_1-6m\,d\bar\zeta_2\right).
\]
Then
\[
 \eta=u\wedge w+6\gamma\wedge\delta,\qquad
 \eta^2=12\gamma\wedge u\wedge w\wedge\delta,\qquad
 Q_F=-d\zeta_1\wedge U-6d\zeta_2\wedge\Gamma.
\]
Because the marked lattice-column determinant is \(y\Delta<0\), the marked
and complex orientations are opposite.  The exact integrals are
\[
 \int_{(F_b,\mathrm{or}_{\rm mark})}Q_F^2=12,\qquad
 \int_{(F_b,\mathrm{or}_{\mathbb C})}Q_F^2=-12,\qquad
 \int_{(F_b,\mathrm{or}_{\mathbb C})}
 Q_F\wedge\Gamma\wedge d\zeta_2=-1.
\]
This is the orientation-correct nonvanishing certificate for the global
section \(\boldsymbol\eta\) and the two Dolbeault inputs.  More precisely,
\(\mathfrak g\) is distinguished and \(\mathfrak q\) is a chosen
coherent-Leray lift, with
\[
 i_b^*\mathfrak g=[\Gamma]_{\bar\partial},\qquad
 i_b^*\mathfrak q=[Q_F]_{\bar\partial}.
\]
The choice of \(\mathfrak q\) is unique only modulo
\(\operatorname{im}d_1\), which restricts trivially to every smooth fibre;
an arbitrary smooth \(Q\) need not restrict coefficient-by-coefficient to
\(Q_F\).

\item Superseded, without deleting it from Round~15, the stale shorthand
\[
 (-C,A,-D,B).
\]
On a finite cover \(t=s^m\), with \((m,a)=(3,2)\) or \((4,1)\), the current
definitions are
\[
 A=\frac{s}{m}d_{00},\quad B=\frac{s^{1-a}}m d_{01},\quad
 C=\frac{s^{1-a}}m d_{10},\quad
 D_0=\frac{s^{m+1-2a}}m d_{11},
\]
and the exact matrices in the actual invariant bases are
\[
 [\kappa_0]=\begin{pmatrix}A&tB\\C&D_0\end{pmatrix},\qquad
 [\kappa_1]=(-C,A,-D_0,B).
\]
The cusp is a separate reduction in its named logarithmic frames:
\[
 [\kappa_0]\equiv\begin{pmatrix}0&-1\\1&0\end{pmatrix}\pmod t,\qquad
 [\kappa_1]\equiv(-1,0,0,-1)\pmod t.
\]
Thus the Round~15 tuple is historical provenance only and is not a current
formula.

\item Completed the associated map-level normalization.  At a finite fibre,
\[
 \operatorname{res}_{1,j}(xb_0+ys^ab_1)=m_jx[b_0],\qquad
 \operatorname{im}(\kappa_0\bmod t)
 =\ker\operatorname{res}_{1,j}=\mathbb C\,s^ab_1.
\]
At the cusp, with
\[
 q_1=(-1,0,1),\quad q_2=(0,-1,-1),\quad
 \alpha_i=\delta_0(q_i),\quad
 \theta=-\delta_1^{-1}(\alpha_1\smile\alpha_2),
\]
the standard \v Cech total differential gives the exact negative sign
\[
 \delta_1\bigl(\rho_*(\alpha\otimes\eta_i)\bigr)
 =-\alpha\smile\alpha_i
\]
and hence the row \((0,1,-1,0)\) in the ordered basis
\((\alpha_1\eta_1,\alpha_1\eta_2,\alpha_2\eta_1,\alpha_2\eta_2)\).
The lift of \(1\) in the evaluation sequence is unique and has
\[
 h=(a(z),1)\binom{d\zeta_1}{d\zeta_2},\qquad a(z_2)=\frac12.
\]
Evaluating in the genuine invariant frame
\[
 \omega_2=s\,u_2(s)(\gamma\wedge u)\ell,\qquad u_2(0)=1,
\]
rather than the frame invariant only at \(s=0\), gives
\[
 \kappa_1(h\otimes\gamma)(p_2)
 =-\frac14\bigl(a(z_2)\dot\tau(z_2)+\dot\mu(z_2)\bigr)
 =-\frac14\dot\mu(z_2)\ne0.
\]

\item Recorded the exact Bott--Chern normalization and its logical scope.
With \(\omega=c_1(\mathcal O_B(1))\),
\[
 \exp(12\mathfrak g)=f^*\mathcal O_B(1),\qquad
 c_{1,BC}(\exp(z\mathfrak g))=z\mathsf A+\bar z\,\mathsf B,
\]
so \(c_{1,BC}(f^*\mathcal O_B(1))=12(\mathsf A+\mathsf B)\).
The algebra includes the unit \(\mathbf1=[1]_{BC}\) and the nine
positive-degree generators span the square-zero ideal \(N\).  Indeed a
\(d\)-closed function on compact connected \(X\) is constant, and there is no
\(\partial\bar\partial\) image in bidegree \((0,0)\); hence
\(H^{0,0}_{BC}(X)=\mathbb C[1]\).  The two
complex-oriented fibre pairings above establish the needed Dolbeault
nonvanishing.  The product \(\mathsf A\mathsf C=0\) uses
Dolbeault--Serre duality together with the independently proved injection
\[
 H^{3,2}_{BC}(X)\hookrightarrow H^{3,2}_{\bar\partial}(X);
\]
the Bott--Chern--Aeppli perfect pairing alone does not force this product to
vanish.  No preferred literal coefficient atlas for the noncanonical smooth
global representatives is claimed.

\item Superseded the interim \(q_{j,\log}\)-deck-group shorthand by the
fully typed affine deck-defect argument (5.17-A14)--(5.17-A15).  The lattice
quotient
\[
 r_j:\widetilde J_j=\Delta_j\times\mathbb C^2\longrightarrow J'_j
\]
has vector-cover automorphisms
\[
 Q_{j,\log}=\widetilde g_j^{\,v_j/m_j},\qquad
 Q_{j,\mathrm{lin}}=\widetilde g_j^{\,0},
\]
with exact relations
\[
 Q_{j,\log}\tau_\lambda Q_{j,\log}^{-1}=\tau_{A_j\lambda},
 \quad (Q_{j,\log})^{m_j}=\tau_{v_j},\qquad
 Q_{j,\mathrm{lin}}\tau_\lambda Q_{j,\mathrm{lin}}^{-1}=\tau_{A_j\lambda},
 \quad (Q_{j,\mathrm{lin}})^{m_j}=1.
\]
The descended automorphisms \(q_{j,\log},q_{j,\mathrm{lin}}\) of \(J'_j\)
have order dividing \(m_j\) because \(\tau_{v_j}\) is trivial in the lattice
quotient, and order exactly \(m_j\) because they cover the exact-order
automorphism \(g_j\).  On \(J_j'^*=J'_j|_{\Delta_j^*}\), the
branch-independent translation \(h_j\) gives
\[
 h_jq_{j,\log}h_j^{-1}=q_{j,\mathrm{lin}}.
\]
The quotient maps and their induced biholomorphism are
\[
\begin{aligned}
 p_{j,\log}:J_j'^*&\to N_j^*,&
 p_{j,\mathrm{lin}}:J_j'^*&\to J^*,\\
 \bar h_j:N_j^*&\xrightarrow{\ \sim\ }J^*,&
 p_{j,\mathrm{lin}}\circ h_j&=\bar h_j\circ p_{j,\log},
\end{aligned}
\]
and the exactly typed line-bundle comparison is
\[
 p_{j,\log}^*(L|_{N_j^*})
 \simeq h_j^*p_{j,\mathrm{lin}}^*(L|_{J^*}).
\]
For the factor extension one returns to the filled vector cover:
\[
 P_{j,\log}:J'_j\to N_j,\qquad
 \pi_j^{\mathrm{vec}}:=P_{j,\log}\circ r_j:\widetilde J_j\to N_j.
\]
Here \(P_{j,\log}\) is the full quotient and
\(P_{j,\log}|_{J_j^{\prime *}}=p_{j,\log}\).  Its deck group is
\[
 \Gamma_j^{\log}
 =\langle\tau_\lambda,Q_{j,\log}:\lambda\in\Lambda\rangle.
\]
The pullback \((\pi_j^{\mathrm{vec}})^*(L|_{N_j})\) is trivial on the
contractible Stein space \(\widetilde J_j\); its factors define an integral
central extension, pushed out along \(\mathbb Z\hookrightarrow\mathbb R\).
Over \(\Delta_j^*\), the line-bundle square identifies each fibre with the
global factor extension by a translation, so its commutator form is \(E\).
The form is integer-valued and locally constant on the connected disc
\(\Delta_j\), hence is \(E\) on the filled cover as well.  Accordingly its
normalized vertical lifts \(U_\lambda^{(j)}\) satisfy
\[
 U_\lambda^{(j)}U_\mu^{(j)}
 =\mathbf c\!\left(\frac12E(\lambda,\mu)\right)
 U_{\lambda+\mu}^{(j)}
\]
For \(\alpha\in\Gamma_j^{\log}\), define \(A_\alpha,T_\alpha\) by
\[
 \alpha\tau_\lambda\alpha^{-1}=\tau_{A_\alpha\lambda},\qquad
 T_\alpha=(A_\alpha^{-1})^{\mathsf t},
\]
and define \(\widetilde d_j(\alpha)\in V_{\mathbb R}\) through
\[
 \widetilde G_\alpha
 U_{A_\alpha^{-1}\lambda}^{(j)}\widetilde G_\alpha^{-1}
 =\mathbf c\!\left(
 \langle\widetilde d_j(\alpha),\lambda\rangle\right)U_\lambda^{(j)},
\qquad
 \widetilde d_j(\alpha\beta)
 =\widetilde d_j(\alpha)+T_\alpha\widetilde d_j(\beta).
\]
The exponent is additive in \(\lambda\), hence determines the unique
\(\widetilde d_j(\alpha)\); composition of lifts proves the crossed law.
For the local-to-global comparison put
\[
 \widetilde J_j^*=\Delta_j^*\times\mathbb C^2,
\]
\[
\begin{aligned}
 \Phi_{j,\log}
  &:=p_{j,\log}\circ r_j|_{\widetilde J_j^*},\\
 \Phi_{j,\mathrm{glob}}
  &:=p_{j,\mathrm{lin}}\circ h_j\circ r_j|_{\widetilde J_j^*}
   =\bar h_j\circ\Phi_{j,\log}.
\end{aligned}
\]
Both maps are \(\Gamma_j^{\log}\)-invariant.  Pulling the line-bundle square
to this common cover gives
\[
 \Theta_j:
 \Phi_{j,\log}^*(L|_{N_j^*})
 \xrightarrow{\ \sim\ }
 \Phi_{j,\mathrm{glob}}^*(L|_{J^*}),
\]
which induces an isomorphism of the pushed-out
\(\Gamma_j^{\log}\)-factor extensions fixing every central
\(\mathbf c(r)\).

Let \(\eta_j\) be the total-space loop obtained by applying the zero section
to the based clockwise meridian \(\rho_j\).  Choose
\(x_j\in J_j^{\prime *}\) over \(\eta_j(0)\), and put
\(y_j=h_j^{-1}(x_j)\).  The \(p_{j,\mathrm{lin}}\)-lift of \(\eta_j\)
from \(x_j\) ends at
\(q_{j,\mathrm{lin}}(x_j)\), while its \(h_j^{-1}\)-transport is the
\(p_{j,\log}\)-lift of \(\bar h_j^{-1}\eta_j\) from \(y_j\), ending at
\(q_{j,\log}(y_j)\).
For a chosen vector lift \(\widetilde y_j\), the adjacent logarithm branch
\[
 \ell'(w)=\ell(g_j^{-1}w)-\frac{2\pi i}{m_j}
\]
has endpoint exactly \(Q_{j,\log}(\widetilde y_j)\), with neither inverse
nor lattice translate.  Therefore, for some \(a_j\in V_{\mathbb R}\) and
\(b_j\in\mathbb R\),
\[
 \Theta_{j,*}(U_\lambda^{(j)})
 =\mathbf c(\langle a_j,\lambda\rangle)U_\lambda,\qquad
 \Theta_{j,*}(\widetilde G_{Q_{j,\log}})
 =\mathbf c(b_j)G_{\rho_j}.
\]
The factor \(\mathbf c(b_j)\) vanishes under conjugation, so
\[
 \widetilde d_j(Q_{j,\log})=d_j+(T_j-I)a_j,\qquad
 \widetilde d_j(\tau_v)=\iota_vE,\qquad
 (\iota_vE)(\lambda)=E(v,\lambda).
\]
Together with \(N_{T_j}(T_j-I)=0\), the vector-cover power relation now
proves, with the positive clockwise sign,
\[
 N_{T_j}d_j
 =\widetilde d_j((Q_{j,\log})^{m_j})
 =\widetilde d_j(\tau_{v_j})=\iota_{v_j}E.
\]
Under \(\ell\mapsto\ell+2\pi in\), one has
\[
 H_j\longmapsto\tau_{nv_j}H_j,\qquad
 \Theta'_{j,*}=\operatorname{Ad}(U_{nv_j})\circ\Theta_{j,*},\qquad
 a_j\longmapsto a_j+n\iota_{v_j}E.
\]
Because \(A_jv_j=v_j\) and \(E\) is \(A_j\)-invariant,
\(T_j(\iota_{v_j}E)=\iota_{v_j}E\).  Hence the correction
\(d_j+(T_j-I)a_j\) and (5.17-A15) are choice-independent.

\item Repaired the final toric restriction presentation.  For adjacent rays
in \(\sigma=\{v,w,u\}\), the actual inclusions are
\[
 U_{\rho_v}=D(z_wz_u)\subset U_\sigma,\qquad
 U_{\rho_w}=D(z_vz_u)\subset U_\sigma,\qquad
 U_{\rho_v}\cap U_{\rho_w}=D(z_vz_wz_u)=U_0.
\]
The contractible Stein chart
\[
 S_\sigma=\{(z_v,z_w,z_u):|z_vz_wz_u|<\epsilon\}
\]
has trivial Picard group and supplies the zero-horizontal-winding comparison.
If \(\mu_\lambda^+\) is multiplication by \(\chi^{m_\lambda}\) on
\(\mathcal O(D_a)\), the negative-divisor map used by the descent tensor
product is the inverse-dual
\[
 \mu_\lambda^-=((\mu_\lambda^+)^{-1})^\vee:
 \Psi_\lambda^*\mathcal O(-D_a)\longrightarrow\mathcal O(-D_a),\qquad
 s\longmapsto\chi^{-m_\lambda}s.
\]
Put \(M=L|_{N_0}\), let \(E_M\) be the Appell--Humbert form of \(M|_{F_t}\),
and fix
\[
 \Lambda_{\rm tor}:=\ker(M_0-I)=\operatorname{im}(M_0-I)
 =\mathbb Z\langle\widehat w,\widehat\delta\rangle,\qquad
 \overline\Lambda:=\Lambda/\Lambda_{\rm tor},
\]
\[
 B_0:\overline\Lambda\xrightarrow{\ \sim\ }\Lambda_{\rm tor},\qquad
 B_0\overline{\widehat\gamma}=-\widehat\delta,\qquad
 B_0\overline{\widehat u}=\widehat w.
\]
Its matrix is
\(\left(\begin{smallmatrix}0&1\\-1&0\end{smallmatrix}\right)\)
from \((\overline{\widehat\gamma},\overline{\widehat u})\) to
\((\widehat w,\widehat\delta)\), not an untyped endomorphism.
For the six degrees
\((d_1,d_2,d_3,d_1,d_2,d_3)\),
\[
 B=\begin{pmatrix}d_2+d_3&d_3\\d_3&d_1+d_3\end{pmatrix},\qquad
 E_M((\lambda,0),(0,m))=-B(B_0\lambda,m)
 \quad(\lambda\in\overline\Lambda,\ m\in\Lambda_{\rm tor}).
\]
The previously proved \(k=0\) gives the indispensable antecedent
\(E_M=c_1(M|_{F_t})=c_1(L|_{F_t})=0\).  Unimodularity of \(B_0\) then
forces \(B=0\), hence
\(d_1=d_2=d_3=0\), and the restriction is trivial.  This round changes
exact presentations and the typed proof-dependency record, not the verified theorem statuses
or their conclusions.  In particular, it explicitly supersedes the
Round~25 symbol \(E_L\) by the typed local bundle \(M=L|_{N_0}\) and its
form \(E_M\), without rewriting that chronological record.
\end{itemize}

\subsection*{Round 27 --- 1 September 2026: independent Hodge replication and primary-source closure}

\begin{itemize}[leftmargin=1.8em]
\item A further context-isolated derivation, performed without treating the
exact-check scripts as proof, independently recovered the two finite cyclic
invariant lattices, the Kodaira--Spencer determinant divisor
$p_1+p_2$, the finite residue-torsion cancellation, the cusp row
$(0,1,-1,0)$, and the order-four scalar
\[
 -\frac14\dot\mu(z_2)\ne0.
\]
It consequently recovered
\[
 R^0f_*\Omega_X^1=\mathcal O(-2),\qquad
 R^1f_*\Omega_X^1=\mathcal O\oplus\mathcal O(-1)^{\oplus2},\qquad
 R^2f_*\Omega_X^1=\mathcal O\oplus\mathcal O(-1),
\]
and hence $h^{1,2}=1$, $h^{1,1}=2$.  This is an independent AI
replication relative to the named family inputs, not an independent human
specialist reconstruction of the entire threefold; the latter remains the
external H2 boundary.
\item The apparent objection $\mathcal E^\vee\not\simeq T_{X/B}$ at the
triple point was tested and refuted.  For
$R=\mathbb C\{x,y,z\}$ and $t=xyz$,
\[
 \ker\!\left[R^3\xrightarrow{(yz,xz,xy)}R\right]
 =R(x,0,-z)\oplus R(0,y,-z)=\mathcal E^\vee.
\]
The three-dimensional pointwise kernel of $df_x$ is not the fibre of this
sheaf kernel: tensoring fails to preserve the kernel because the image ideal
$(xy,xz,yz)$ is not flat.  The workbench now states this distinction
explicitly.
\item The actual primary sources were downloaded into the private local
literature corpus and read at the relevant theorem pages.  Atiyah--Singer,
Theorem~3, p.~427 proves HRR for every compact complex manifold and
holomorphic vector bundle; Serre, \S13, Theorem~4, p.~22 gives the compact
complex duality used in the Hodge table; and the public Angella--Tomassini
LaTeX source proves $h^k_{BC}+h^k_A\ge2b_k$.  All three are now cited
directly.  The literature ledger now limits Cartan, Deligne, and Katz--Oda
to their exact roles instead of attributing the candidate's singular-fibre
maps to smooth-background references.
\item A bounded provenance search found no authenticated original
Claude/Fabel session export, model-session identifier, or author-controlled
model record.  The degree-seven polynomial collision certificate remains
independently verified, but its discovery attribution remains
\textsc{open}; no mathematical claim status changes.
\end{itemize}

\subsection*{Round 28 --- 3 September 2026: cusp-cover scope and T1/T2 exact-morphism audit}

\begin{itemize}[leftmargin=1.8em]
\item The native radial family and the later finite-cover construction are now
separated at the level of exact objects.  Projection onto the two imaginary
coordinates gives
\[
 0\longrightarrow\mathbb T^2_{\rm re}\longrightarrow F_T
 \longrightarrow\mathbb R^2/Y_T\mathbb Z^2\longrightarrow0,
 \qquad
 Y_T=\begin{pmatrix}6m_T&L_T\\-q_T&m_T\end{pmatrix},
 \qquad \det Y_T=D_T,
\]
so the unmodified S6 regular-fibre family has native pointed limit
\(\mathbb T^2\times\mathbb R^2\).  The separate sublattices
\(\mathsf S_M\mathbb Z^4\), \(\mathsf S_M=\operatorname{diag}(1,1,M,M)\),
give degree-\(M^2\) covers only after the base change \(t_c=v^M\), with
\[
 \mathsf S_M^{-1}M_0^M\mathsf S_M=M_0.
\]
Their pointed \(\mathbb R^4\) limit is valid but auxiliary: no extension over
the central fibre, global cover of \(X\), or S12/S24-style global
generalization is thereby constructed.

\item The T1 spectral packet contained one S6-specific theorem not already in
the workbench.  It has now been reconstructed from the verified canonical
period functions.  The first coordinate induces
\[
 \mathcal A_z\longrightarrow\mathbb C/(\mathbb Z\tau+\mathbb Z)
\]
exactly when \(6\mu=a\tau+b\) for a unique \((a,b)\in\mathbb Z^2\).  Its
connected kernel is
\[
 \mathbb C/(\mathbb Z+\mathbb Z(a\mu-\beta)),
 \qquad \operatorname{Im}(a\mu-\beta)=-D_{\rm can}>0.
\]
The marked labels transform by
\[
 g_0(a,b)=(a,a+b),\qquad g_2(a,b)=(-b,6+a).
\]
Allowing the labels to vary, the identity theorem and the bounded cusp
extension of \(\mu\) exclude every cusp-accumulating sequence.  Compactness
of a truncated triangle fundamental polygon then proves that the locus has
finite image in the completed orbifold base.  This classifies only this
specified projection; no claim about every elliptic quotient or integral
magnetic class is made.

\item The independent T1 Koide arithmetic was audited but not inserted into
the \(S^6\) proof.  The audit covered the Vieta descent, the exact bijection
\[
 (x,y,z)\longleftrightarrow(a,q,n),\qquad
 a^2-3q^2-18n^2=-2,\qquad z=4a+17n,
\]
the local densities, and the archimedean factor check.  At the Round~28
cutoff, the claimed linear counting constant still required an explicit
verification of the Gorodnik--Oh--Shah sector hypotheses and a
measure-normalization comparison with Borovoi; Round~29 records the completed
proof of both points.  The displayed 2-adic element is a sign-reversing
bijection but not an involution.  Since no morphism from a regular
abelian-surface fibre of the \(S^6\) family to
the defect-one integral shell is supplied, this arithmetic remains a separate
preprint line, cross-referenced only through the proved ambient determinant
identity.
\end{itemize}

\subsection*{Round 29 --- 3 September 2026: exact four-copy globalization,
cusp obstruction, and Koide normalization closure}

\begin{itemize}[leftmargin=1.8em]
\item The regular period family was rechecked at the level of its group law.
Its linear monodromy fixes the zero section, so
\(J=X^\circ\to B^\circ\) is a holomorphic family of complex Lie groups,
not merely an affine torus bundle.  With
\[
 K_\Lambda=\mathbb Z\langle\widehat w,\widehat\delta\rangle,
 \qquad
 L_\Lambda=\mathbb Z\langle\widehat u,
                    \widehat w,\widehat\delta\rangle,
\]
\(L_\Lambda\) is the smallest full-monodromy-invariant lattice containing
\(K_\Lambda\), both \(K_\Lambda\) and \(L_\Lambda\) are primitive, and
\[
 \operatorname{Hom}(\Lambda,\mathbb Z)^\Delta
 =\mathbb Z\varepsilon_\gamma.
\]
The primitive character \(\varepsilon_\gamma\) gives the global real group
quotient \(J\to B^\circ\times S^1\) with connected three-torus kernel.  Its
radial quotient circumference is exactly
\[
 H_{\mathrm{cusp}}:=\operatorname{Im}h(0),\qquad
 B_{\mathrm{cusp}}:=\operatorname{Im}b(0),
 \qquad
 \ell_T=\frac{D_T}{\sqrt{L_T^2+m_T^2}}
       =T+H_{\mathrm{cusp}}-B_{\mathrm{cusp}}+O(T^{-1}).
\]
Four synchronized copies therefore give a genuine \(\mathbb T^4\)-quotient
whose pointed radial limit is \(\mathbb R^4\); the invariant-lattice rank
argument proves that four copies are minimal.

\item The tempting two-copy route cannot be globalized by passing to a finite
cover of the base.  For \(H_*=A_2A_1=I+N_*\), one has \(N_*^2=0\) and
\[
 H_*^n\widehat w
 =\widehat w+n(\widehat u-\widehat\delta),
\]
so no positive power preserves \(K_\Lambda\).  At the two finite fillings,
the character instead extends equivariantly to explicit Seifert targets;
the fourfold simultaneous-root models normalize the ordinary fibre products
with reduced-central-fibre degrees \(3^3=27\) and \(4^3=64\).

\item The literal quotient does not extend through the toric cusp.  Two
explicit radial paths approach the same triple point of \(W\), while their
character values remain \(3/4\) and \(1/2\).  This proves continuous
nonextension of both the one-copy character and its fourfold product.  The
topological alternatives are also typed exactly:
\[
 X^2\simeq S^6\times S^6\ne S^{12},\qquad
 X^4\simeq(S^6)^4\ne S^{24},
\]
whereas \(X\wedge X\cong S^{12}\) and
\(X^{\wedge4}\cong S^{24}\), but the quotient maps are nonconstant on the
relevant fat wedges.  Borel--Serre, \S15, Proposition~15.1 independently
excludes almost-complex structures on \(S^{12}\) and \(S^{24}\), for every
differentiable structure.  At the Round~29 cutoff a different compact or
singular globalization was still open.  Round~31 constructs such
globalizations by replacing the fixed target with cone targets; the
displayed literal fixed-target extension remains refuted.

\item The Koide defect-one counting normalization left open in Round~28 is
now closed.  The exact Vieta bijection identifies reduced positive roots with
\[
 a^2-3q^2-18n^2=-2,\qquad
 a>0, q,n\ge0, a+4n>q,\qquad z=4a+17n.
\]
The bad-prime densities, the good-prime Euler factors, the global
Poincar\'e-residue measure, the real chamber integral, the
Gorodnik--Oh--Shah admissibility and genericity conditions, and Borovoi's
adelic sign coefficient have all been proved in the displayed coordinates.
The resulting unconditional theorem is
\[
 R(X)\sim
 \frac9{2\pi^2}\log(2+\sqrt3)
 \log\!\left(\frac3{2\sqrt2}\right)X.
\]
Only the sharper
\(O_\varepsilon(X^{1/2+\varepsilon})\) discrepancy remains open.

\item The scalar, cover, continuum gauge, magnetic, and Jacobi operators were
also given their explicit \(L^2\) Hilbert spaces, closed forms, and
self-adjoint realizations.  In particular, the full background-gauge Jacobi
operator is a connection Laplacian plus a bounded symmetric endomorphism on
\(H^2\), its form domain is \(H^1\), and the Coulomb slice reduces it by the
proved resolvent intertwining with the scalar covariant Laplacian.  This
closes the operator-domain verification without changing the boundary between
the verified classical spectra and the open physical Yang--Mills transfer.
\end{itemize}

\subsection*{Round 30 --- 3 September 2026: exact
Erd\H{o}s--Straus/S6 relation-matrix crosswalk}

\begin{itemize}[leftmargin=1.8em]
\item The complete first-half divisor-shell selector for a prime
\(p\equiv1\pmod{12}\) was derived again from
\((R_jy-pa_j)(R_jz-pa_j)=p^2a_j^2\), including the historically named
exponent-zero/two complement channel and the exponent-one middle channel.
In the selector below the channel variable itself satisfies
\(\varepsilon\in\{0,1\}\).  No universal
nonemptiness statement is used.

\item With \(R_j=4j+3\), a selector state \(u\mid(a_j^2)\),
\(\varepsilon\in\{0,1\}\), and
\(R_jT=4u+p^\varepsilon\), the exact coordinate change
\[
 \ell_0=\frac{T+1}{4},\qquad \ell_1=u-jT,\qquad \ell_2=1
\]
gives \(12\ell_0-4\ell_1-3=p^\varepsilon\) and the unimodular identity
\[
 \begin{pmatrix}1&-j\\0&1\end{pmatrix}
 \begin{pmatrix}R_j&-u\\4&-T\end{pmatrix}
 =\begin{pmatrix}3&-\ell_1\\4&1-4\ell_0\end{pmatrix}.
\]
The inverse is \(T=4\ell_0-1\), \(u=\ell_1+jT\).  Retaining the prime,
shell mark, channel, orientation, and divisor condition makes the map
bijective.  The result is exactly an integer-presentation equivalence; it
does not manufacture a morphism between the reciprocal cubic and the
candidate threefold.

\item A finite independent replay over all (300) primes
\(p<10{,}000\), \(p\equiv1\pmod{12}\), checked (358{,}959) shells,
\(21{,}084\) selector states, and (32{,}576) raw divisor states: every
integrality, matrix, determinant, inverse, denominator, and
reciprocal-equation identity passed exactly, with zero errors.
\end{itemize}

\subsection*{Round 31 --- 3 September 2026: local Kummer completion,
cone-smash sphere carriers, and compact-link/Gaussian audit closure}

\begin{itemize}[leftmargin=1.8em]
\item The punctured degree-$N^2$ fibre cover has now been extended through
the toric cusp, but only after replacing the source action by the exact
base-changed action.  With $t=v^N$ one must use
\[
 C_N(v)=\frac1N C(v^N),\qquad
 \Psi^{(N)}_{\bar\lambda}(\xi,v)
 =\bigl(\exp(2\pi iC_N(v)\bar\lambda)
         v^{B_0\bar\lambda}\xi,v\bigr),
\]
not the unmodified target action.  Multiplication by $N$ on the three
toric cocharacter directions gives
\[
 P_N(\xi_1,\xi_2,v)=(\xi_1^N,\xi_2^N,v^N),
\]
a finite flat map of degree $N^3$, equivariant for the modified source
action and the original target action.  It descends to
\[
 W_N:=\bigl(f_N^{-1}(0)\bigr)_{\mathrm{red}},\qquad
 \kappa_N:N_0^{\langle N\rangle}\longrightarrow N_0,
 \qquad f_0\circ\kappa_N=\rho_N\circ f_N,
\]
of degree $N^3$, with $\kappa_N^*W=NW_N$ and
$\operatorname{Ram}(\kappa_N)=(N-1)W_N$.  Factoring through the normal base
change $N_0[N]=N_0\times_{\Delta_t}\Delta_v$ gives the relative map
\[
 N_0^{\langle N\rangle}\xrightarrow{Q_N}N_0[N]
\]
of generic degree $N^2$.  For $N>1$, $Q_N$ is quasi-etale but not flat:
at a triple stratum its fibre length is
$N^3-(N-1)^3=3N^2-3N+1>N^2$, and its non-etale and nonflat loci are the
inverse images of the double and triple strata.  On reduced central fibres
the generic degree is $N^2$; on their normalizations the map is toric
multiplication $[N]:\mathrm{dP}_6\to\mathrm{dP}_6$, of degree $N^2$ and
ramification index $N$ along each of the six boundary curves.  Hence the
completed local map to the original family has degree $N^3$, while degree
$N^2$ belongs to the relative base-changed map and the reduced-central-fibre
map.  This is not a degree-$N^2$ self-map over the unchanged disc, not the
normalization of the base change, and not a cover over the whole base $B$;
by itself it supplies no $S^{12}$ or $S^{24}$ map.

\item Three different cone-smash routes have now been constructed while
keeping their sources, targets, and collapse maps distinct.  The native
punctured-cusp quotient modulo
$K_\Lambda=\mathbb Z\langle\widehat w,\widehat\delta\rangle$ gives
\[
 q_2:(X,x_0)\longrightarrow
 \mathcal A=\operatorname{Susp}(S^1\times\overline T^2)
             \cong\operatorname{Susp}(T^3),
\]
where $W$ and $X\setminus\operatorname{int}N_0$ go to the two suspension
vertices.  Its smash square
\[
 q_2^{\wedge2}:X\wedge X\cong S^{12}
 \longrightarrow\mathcal A\wedge\mathcal A
\]
restricts on the embedded radial slice $F_T^2$ to the literal product
$\pi_T\times\pi_T:F_T^2\to Q_T^2$, whose target tends intrinsically to
$\mathbb R^4$.  The target is not a sphere or compact four-manifold:
$\widetilde H_*(\mathcal A)$ has ranks $3,3,1$ in degrees $2,3,4$, and
$\mathcal A^{\wedge2}$ has ranks $9,18,15,6,1$ in degrees $4,5,6,7,8$.
In particular, it is $\mathcal A^{\wedge2}$---not the compact four-torus
$Q_T^{\times2}$---that is neither a sphere nor a compact four-manifold.

The cusp-local invariant character separately gives
\[
 q_{1,\mathrm{cusp}}:(X,x_0)\longrightarrow
 \mathcal B=\operatorname{Susp}(T^2),
 \qquad
 q_{1,\mathrm{cusp}}^{\wedge4}:S^{24}
 \longrightarrow\mathcal B^{\wedge4},
\]
whose synchronized radial restriction is
$\chi_{\gamma,T}^{\times4}:F_T^4\to C_T^4$.  Here
$\widetilde H_*(\mathcal B)$ has ranks $2,1$ in degrees $2,3$, and the
fourth smash has ranks $16,32,24,8,1$ in degrees $8,9,10,11,12$.
For the full-base route, gluing the two finite-filling Seifert targets to
$P\times S^1$ gives a compact orientable three-manifold $\mathcal M$ with
boundary $T^2$.  Coning the entire boundary gives
\[
 q_1:(X,x_0)\longrightarrow
 \mathcal C=\mathcal M\cup_{\partial\mathcal M}
             \operatorname{Cone}(\partial\mathcal M),
 \qquad
 q_1^{\wedge4}:S^{24}\longrightarrow\mathcal C^{\wedge4}.
\]
Its radial restriction is again $\chi_{\gamma,T}^{\times4}$, with
$\ell_T=T+H_{\mathrm{cusp}}-B_{\mathrm{cusp}}+O(T^{-1})$ and intrinsic
pointed limit $\mathbb R^4$.
One has $\widetilde H_k(\mathcal C)=\mathbb Z$ for $k=2,3$ and zero
otherwise, while $\mathcal C^{\wedge4}$ has ranks $1,4,6,4,1$ in degrees
$8,9,10,11,12$.  These global maps exist because the codomains retain the
radial parameter and collapse the obstructing boundary cross-section.
They do not extend the fixed-target torus quotient, make $X^2$ or $X^4$
into spheres, or contradict the fat-wedge nonfactorization: the sphere
sources are the based smash spaces $X\wedge X\cong S^{12}$ and
$X^{\wedge4}\cong S^{24}$, and the radial products embed because the
chosen global base point lies in $W$.  The limits are intrinsic flat-slice
limits; in any fixed compatible metric on a compact cone target the same
cross-sections collapse toward the cone vertex.  Neither $S^{12}$ nor
$S^{24}$ admits an almost-complex structure, so these real carriers are
not compact-complex sphere globalizations.

\item The compact-link and Gaussian calculations were audited after
separating the cover index $M_{\mathrm{cov}}\geq1$ from the mesh size
$N_{\mathrm{mesh}}\geq2$.  The regular-fibre degree-$M_{\mathrm{cov}}^2$
cover has
\[
 P_{T,M_{\mathrm{cov}}}
 =P_T\operatorname{diag}(1,1,M_{\mathrm{cov}},M_{\mathrm{cov}}),
 \qquad \operatorname{covol}=M_{\mathrm{cov}}^2D_T,
\]
whereas the transported cubical subdivision has
$N_{\mathrm{mesh}}^4$ vertices.  The displayed vertex map reaches exactly
the finite $N_{\mathrm{mesh}}$-torsion subgroup, not the whole torus.
For the explicit unit-flux configuration,
\[
 H_{\mathrm{link}}:=\operatorname{diag}(i,-i),\qquad
 U_{12}(n)=\exp(2\pi H_{\mathrm{link}}/N_{\mathrm{mesh}}^2)
\]
also across both wrap boundaries; all other plaquettes with $i<j$ are
trivial, while $U_{21}=U_{12}^{-1}$ is not an additional trivial
plaquette.  The logarithmic action and the explicitly defined
continuum-matched $(1,2)$-sector Wilson action are
\[
 S^{\log}=\frac{2\pi^2M_{\mathrm{cov}}^2}
                  {g_{\mathrm{YM}}^2D_T},\qquad
 S^{W,12}=\frac{M_{\mathrm{cov}}^2N_{\mathrm{mesh}}^4}
                  {2g_{\mathrm{YM}}^2D_T}
 \left(2-2\cos\frac{2\pi}{N_{\mathrm{mesh}}^2}\right).
\]
The Gaussian Hamiltonian is the Friedrichs oscillator on the finite
noncompact linear Coulomb slice, with gap
$\sqrt{\lambda_{\min}(g_{\mathrm{YM}}^{-2}d^*d)}$, and the soft mode gives
\[
 \operatorname{gap}\le
 \frac{2N_{\mathrm{mesh}}\sin(\pi/N_{\mathrm{mesh}})}{g_{\mathrm{YM}}}
 \frac{\sqrt{L_T^2+m_T^2}}{D_T}.
\]
The exact simultaneous sequence
\[
 T_j=j^2,\qquad M_{\mathrm{cov},j}=j,\qquad
 N_{\mathrm{mesh},j}=j^3
\]
gives systole at least
$\min\{T_j-C,M_{\mathrm{cov},j}\}\to\infty$, maximum physical mesh edge
$O(j^{-1})$, Gaussian-gap upper bound $O(j^{-2})$, and logarithmic action
$O(j^{-2})$.  The audit also fixes the physical boundary: the componentwise
exponential has derivative $N_{\mathrm{mesh}}^{-1}$, is only a local chart,
and is not globally injective.  If $DS_W(\mathbf1)=0$, then the exact
critical-point chain rule gives
$D^2(S_W\circ\operatorname{Exp}_{N_{\mathrm{mesh}}})_0=
N_{\mathrm{mesh}}^{-2}D^2S_W(\mathbf1)$ under coefficient identification.
The Gaussian Lebesgue $L^2$-space is not the
compact-link Haar $L^2$-space; no unitary identification, nonlinear
compact-link Hamiltonian, transfer matrix, renormalized limit, or
Osterwalder--Schrader reconstruction is supplied, and a link configuration
defines a Dirac distribution rather than an $L^2$ vector.

\item The independent receipts are recorded with paths and exact byte scope.
The temporary snapshot
\texttt{tmp/agent\_results/cone\_compactified\_sphere\_route.tex}
has 30,662 bytes and SHA-256
\texttt{F09F133115887998FA4C7510E71EABCC5B7E68F0135C0F673E7BFB5D11E82F31};
its clean audited wrapper PDF
\texttt{tmp/agent\_results/cone\_compactified\_sphere\_route\_audited\_wrapper.pdf}
has 237,982 bytes and SHA-256
\texttt{6767B5CDF5788D3A891B33968C2304DF94E6C5416315C3895DA4CA1D1BAF1E0F}.
These are explicitly historical temporary snapshots, not claims about the
current canonical source.  The accompanying proof audit closed the earlier
integral-homology gap: Euler characteristic alone gave only a rank statement
for $H_2(\mathcal M)$; Poincar\'e--Lefschetz duality and the universal
coefficient theorem prove this group is free, so rank zero gives
$H_2(\mathcal M)=0$.  It also supplied the closed-fat-wedge and Hausdorff
arguments needed for the slice embeddings.

The file
\texttt{tmp/agent\_results/compact\_link\_symbolic\_receipt.md}
has 4,151 bytes and SHA-256
\texttt{1C7BC6F3BD3FE0B769E43304DB826FCE70241206C797572220BEFF2745011E1D}.
It records 17,214 exact wrap/plaquette assertions together with independent
Gram-minor, action, Fourier-symbol, soft-mode, and asymptotic checks, with no
discrepancy.  The earlier audited-input digest printed inside that receipt is
a temporary-snapshot identifier; the corresponding exact input bytes are not
presently retained under a verified path, so that digest is not equated with
the canonical compact-link source.  Likewise, the earlier unsupported claim
that a Kummer staging file was byte-identical to the canonical source is
withdrawn: no separately retained staging path establishes it.

For reproducibility only, a non-final checksum checkpoint was taken at
2026-09-03 05:17:48 +02:00 while propagation was still active.  At that
instant the following canonical-root-relative paths had the displayed byte
counts and SHA-256 digests:
\begin{itemize}[leftmargin=2em]
\item \texttt{supporting\_materials/workbench/research/
cusp\_cone\_sphere\_routes.tex}: 31,893 bytes;\\
\texttt{3CD16F10A13EC6F740A67A86E48F74E1880438D10F1CE915CB15D3267D2FB26D}.
\item \texttt{supporting\_materials/workbench/research/
cusp\_sphere\_route\_strengthening.tex}: 28,220 bytes;\\
\texttt{B397722C2BCC8429AD4957EA51C6D70ADE4ABB1651CE3EF988A2CD38942D8623}.
\item \texttt{supporting\_materials/workbench/research/
cusp\_sphere\_fixed\_metric.tex}: 17,586 bytes;\\
\texttt{3072076F55A8DCD32787C86DB03A0523F5B8FA7D80C0F54704791607C0533C0B}.
\item \texttt{supporting\_materials/workbench/research/
cusp\_compact\_link\_gaussian.tex}: 32,214 bytes;\\
\texttt{4469781B2386225453DBDD5C12B5B5D9213CD67480055F3F318143941C44FA52}.
\item \texttt{supporting\_materials/workbench/research/
cusp\_toric\_cover\_completion.tex}: 16,474 bytes;\\
\texttt{076A6CC0D6F8BD9E86282C6C088B867FBCB2FC8683F3550CE0896EB04CDD3701}.
\item \texttt{s6\_spectral\_degeneration\_preprint.tex}: 62,448 bytes;\\
\texttt{A83FAF8093A85712033BFADCAAB81E8408023D8883B6DE8D6AB95C826258ACB2}.
\end{itemize}
Because this was a propagation checkpoint rather than a frozen release, it
is not described as current; final release hashes must be regenerated after
the canonical tree freezes.
\end{itemize}

\subsection*{Round 32 --- 3 September 2026: native sphere carriers,
moving four-spheres, and exact low-energy populations}

\begin{itemize}[leftmargin=1.8em]
\item The finite-cover construction and the requested sphere-carrier problem
are now separated explicitly.  The former is a correct ancillary theorem:
for
\[
 S_N=\operatorname{diag}(1,1,N,N),\qquad
 \widetilde F_{T,N}=\mathbb R^4/P_TS_N\mathbb Z^4,
\]
the identity on \(\mathbb R^4\) induces an \(N^2\)-sheeted holomorphic
local-isometric cover of \(F_T\), and after \(t_c=s^N\) one has
\[
 S_N^{-1}M_0^NS_N=M_0,
 \qquad
 \operatorname{sys}(\widetilde F_{T,N})\geq\min\{T-C,N\}.
\]
Thus \(T_j,N_j\to\infty\) gives a pointed \(\mathbb R^4\) limit and the
cover characters fill the two momentum directions absent from pullback of
the original torus modes.  This changes the fibre lattice and the cusp base.
It therefore does not answer the distinct question of deriving an
\(S^{12}\)- or \(S^{24}\)-carried \(\mathbb R^4\) limit from the native
\(\mathbb T^2\times\mathbb R^2\) quotient.  No cover of \(X\) or of the
full base is asserted.

\item The native sphere-carrier route has now been strengthened directly.
After fixing a based diffeomorphism \(X\cong_{\mathrm{diff}}S^6\), there
are continuous surjective based quotient maps
\[
 \widehat q_3^{\wedge2}:S^{12}\twoheadrightarrow S^6,
 \qquad
 \widehat q_\gamma^{\wedge4}:S^{24}\twoheadrightarrow S^8.
\]
The first target contains the moving based four-sphere
\(\mathbb S_{12,T}^4=S_T^2\wedge S_T^2\), obtained from
\(F_T^{\times2}\to Q_T^{\times2}\) by collapsing the coordinate fat
wedge.  The second contains
\(\mathbb S_{24,T}^4=C_T^{\wedge4}\), obtained from
\(\chi_{\gamma,T}^{\times4}:F_T^{\times4}\to C_T^{\times4}\) by
collapsing the character fat wedge.  In their quotient length metrics,
the midpoint radius-\(R\) ball is exactly Euclidean when
\(T-C>4R\), respectively when \(\ell_T>4R\).  Hence both moving
four-spheres converge to \(\mathbb R^4\) pointed locally isometrically,
smoothly, and in pointed Gromov--Hausdorff topology.  Sliding the reduced
suspension coordinate contracts both inclusions, so these four-spheres are
null-homotopic in \(S^6\) and \(S^8\) and do not contradict the fixed-target
orientation obstruction.  The unscaled quotient-norm rank calculation proves
that two copies are minimal for the native cusp-local rank-four quotient and
four copies are minimal for a rank-four quotient surviving full-base
monodromy.

\item The target metric scope is also exact.  For every fixed compatible
metric on any compact cone target, the synchronized radial slices converge
in Hausdorff and Gromov--Hausdorff topology to the smash basepoint.  There
are nevertheless fixed cusp-adapted compatible length metrics for which
moving-point blow-ups recover the intrinsic \(\mathbb R^4\) limit on the
vertical slices.  Under the same rescaling the full ambient bubbles are
\(\mathbb R^8\) for the \(S^{12}\) route and \(\mathbb R^{12}\) for both
\(S^{24}\) routes.  If \(\beta_T=\lambda_Tr_T\), the complete scalar
rescaling trichotomy for a slice is
\[
 \beta_T\to0:\{*\},\qquad
 \beta_T\to L\in(0,\infty):
 (\mathbb R/L\mathbb Z)^4,\qquad
 \beta_T\to\infty:\mathbb R^4.
\]
The unit-volume normalization gives the square flat four-torus.  Every
rescaling which unrolls the slice sends the cone vertex to infinite distance,
so the \(\mathbb R^4\) limit is a moving slice bubble, not a vertex tangent
cone or a compact-target asymptotic cone.  The quotient construction supplies
no Riemannian metric on the source spheres.

\item The T1 scalar spectral packet contributes a fixed-energy theorem not
contained in the earlier first-gap statement.  With
\[
 \mathcal D_T=|\det P_T|=T^2(1+O(T^{-1}))
\]
and Fourier multiplicities retained, every fixed \(E\geq0\) satisfies
\[
 N_T(E)=\pi\mathcal D_T
 \sum_{n\in\mathbb Z^2}
 \left(\frac{E}{4\pi^2}-\lVert n\rVert^2\right)_+
 +O_E(T).
\]
For \(0<E<4\pi^2\) only \(n=0\) occurs, and therefore
\[
 N_T(E)=\frac{E\mathcal D_T}{4\pi}+O_E(T)
       =\frac{ET^2}{4\pi}+O_E(T).
\]
The proof solves the complete inverse-transpose lattice equation, reduces
each transverse band to a translated planar ellipse, and uses an exact
area--perimeter lattice estimate.  Thus the low-energy population has a
two-dimensional leading density; this is not a Weyl law for a fixed metric
on the compact total space.

\item The T2 magnetic calculation gives an independent full-torus result.
Every nondegenerate integral alternating form
\(\mathsf E\in\bigwedge^2(\mathbb Z^4)^*\) defines the explicit factor of
automorphy line bundle with
\[
 \nabla^{\mathsf E}=d+2\pi i\sum_{j<k}E_{jk}u_j\,du_k,
 \qquad F_{\nabla^{\mathsf E}}=2\pi i\mathsf E.
\]
The closed magnetic form has a Friedrichs operator of domain \(H^2\) and
compact resolvent.  For the two positive skew frequencies \(b_1,b_2\), its
complete set of spectral values is
\[
 \{(2n_1+1)b_1+(2n_2+1)b_2:n_1,n_2\in\mathbb Z_{\geq0}\},
\]
and
\[
 b_1b_2=\frac{4\pi^2|\operatorname{Pf}(\mathsf E)|}{\mathcal D_T},
 \qquad
 \operatorname{gap}\leq
 4\pi\sqrt{\frac{|\operatorname{Pf}(\mathsf E)|}{\mathcal D_T}}
 =O(T^{-1}).
\]
This construction uses the full marked four-torus and does not require the
exceptional first-coordinate elliptic quotient.  It does not assert a
full-monodromy-invariant global flux class and is not a Yang--Mills mass-gap
theorem.

\item The claim ledger now includes the previously unlisted open Koide
discrepancy.  The unconditional theorem
\(R_{\mathrm K}(Y)\sim C_0Y\) remains verified; only
\(R_{\mathrm K}(Y)-C_0Y=O_\varepsilon(Y^{1/2+\varepsilon})\) is open.  After
the Round~32 additions the machine-readable claim ledger contains 64
\texttt{VERIFIED}, 4 \texttt{OPEN}, 5 \texttt{REFUTED}, and no
\texttt{CONDITIONAL} rows.
\end{itemize}

\subsection*{Round 33 --- 3 September 2026: correction of the
boundary-family comparison topology}

\begin{itemize}[leftmargin=1.8em]
\item The comparison maps from the exact decompactification boundary
families to the compact cone carriers have been re-audited.  For
\(\nu\in\{12,24\}\),
\[
 \beta_\nu:\mathfrak E_\nu\longrightarrow\mathcal D_\nu
\]
is continuous and surjective, is a homeomorphism away from the zero fibre,
and satisfies
\[
 \beta_\nu^{-1}(d_\nu)=p_\nu^{-1}(0)\cong\mathbb R^4.
\]
It is nonproper because the inverse image of the compact singleton
\(\{d_\nu\}\) is the noncompact space \(\mathbb R^4\).

\item The stronger quotient interpretation is false.  Fix
\(a\in\mathbb R^4\setminus\mathbb Z^4\), choose \(s_j\downarrow0\), and put
\[
 z_j=[s_j,A_{1/s_j}^{(\nu)}a]\in\mathfrak E_\nu.
\]
For \(\nu=12\),
\[
 \inf_{n\in\mathbb Z^4}
 \|A_{1/s_j}^{(12)}(a+n)\|
 \geq(1/s_j-C)\operatorname{dist}(a,\mathbb Z^4)\longrightarrow\infty,
\]
whereas for \(\nu=24\) the same infimum equals
\[
 \ell_{1/s_j}\operatorname{dist}(a,\mathbb Z^4)\longrightarrow\infty.
\]
Hence \(Z=\{z_j:j\geq1\}\) is closed in \(\mathfrak E_\nu\): it has no
positive-fibre accumulation point and eventually leaves every half-space
chart at the zero fibre.  The set is \(\beta_\nu\)-saturated because
\(\beta_\nu\) is injective away from the zero fibre and \(\beta_\nu(Z)\)
does not contain the vertex.  On the other hand,
\[
 \beta_\nu(z_j)=[e^{-2\pi/s_j},a]\longrightarrow d_\nu,
\]
so \(\beta_\nu(Z)\) is not closed.  Therefore \(\beta_\nu\) is not a
quotient map and does not identify \(\mathcal D_\nu\) with
\(\mathfrak E_\nu/p_\nu^{-1}(0)\).  In particular it is not a topological
blow-down or resolution.

\item This correction does not alter the exact positive content:
\(\mathfrak E_\nu\) remains a Hausdorff second-countable topological
five-manifold with boundary \(\mathbb R^4\), its positive fibres remain the
native expanding four-tori, and their pointed local-isometric, smooth, and
Gromov--Hausdorff convergence to the boundary fibre remains verified.  The
moving four-spheres, their intrinsic \(\mathbb R^4\) limits, and the
fixed-metric collapse and scaling trichotomy are independent of the false
quotient-space interpretation.

\item After adding the explicit refutation, the claim ledger contains 64
\texttt{VERIFIED}, 4 \texttt{OPEN}, 6 \texttt{REFUTED}, and no
\texttt{CONDITIONAL} rows; the morphism ledger contains 49
\texttt{VERIFIED}, 1 \texttt{OPEN}, and 2 \texttt{REFUTED} rows; the
uncertainty ledger contains 30 \texttt{VERIFIED}, 5 \texttt{OPEN}, and 4
\texttt{REFUTED} rows; and the dependency ledger contains 152 rows.
\end{itemize}

\subsection*{Round 34 --- 3 September 2026: literal native
\(\mathbb R^4\) cells inside the sphere targets}

\begin{itemize}[leftmargin=1.8em]
\item The topological content of the moving-sphere theorem is now stated
without passing to a limit.  For every radial parameter \(T\), set
\[
 U_{12,T}=\mathbb S_{12,T}^4\setminus\{*\},
 \qquad
 U_{24,T}=\mathbb S_{24,T}^4\setminus\{*\}.
\]
Restriction of a closed-subset collapse to the complement of its quotient
point is a homeomorphism.  Since the coordinate one-skeleton complement in
the marked two-torus is its unique open two-cell and a based-circle
complement is an open interval, the exact quotient-stage restrictions are
\[
 (Q_T\setminus K_T^{(1)})^2
 \xrightarrow{\ \cong\ }U_{12,T}\cong\mathbb R^4,
\]
\[
 (C_T\setminus\{0\})^4
 \xrightarrow{\ \cong\ }U_{24,T}\cong\mathbb R^4.
\]
Composing with the already verified moving-four-sphere embeddings gives
literal topological embeddings
\[
 U_{12,T}\cong\mathbb R^4\hookrightarrow S^6,
 \qquad
 U_{24,T}\cong\mathbb R^4\hookrightarrow S^8.
\]

\item The word ``open'' has its exact relative meaning.  Each cell is open
in its moving four-sphere, is not open in the ambient \(S^6\) or \(S^8\),
and is nonclosed there; its ambient closure is the corresponding moving
four-sphere.  It is not globally isometric to Euclidean \(\mathbb R^4\)
under either the assigned compact quotient metric or the restriction of a
fixed compatible compact-target metric.  The precise metric assertion is
instead the already proved exhaustion by exact midpoint balls: for every
fixed \(R\), the pointed radius-\(R\) ball is Euclidean and the isometry is
onto once \(T-C>4R\), respectively once \(\ell_T>4R\).

\item This is a native quotient-stage theorem and not a reformulation of the
ancillary finite-cover construction.  It changes neither the fibre lattice
nor the cusp base and introduces no cover modes.  It is also independent of
the nonquotient boundary-family comparison \(\beta_\nu\): the cell maps are
restrictions of the genuine compact-to-Hausdorff smash quotient maps.

\item After this propagation, the claim ledger contains 65
\texttt{VERIFIED}, 4 \texttt{OPEN}, 6 \texttt{REFUTED}, and no
\texttt{CONDITIONAL} rows; the morphism ledger contains 51
\texttt{VERIFIED}, 1 \texttt{OPEN}, and 2 \texttt{REFUTED} rows; the
uncertainty ledger contains 31 \texttt{VERIFIED}, 5 \texttt{OPEN}, and 4
\texttt{REFUTED} rows; and the dependency ledger contains 157 rows.
\end{itemize}

\subsection*{Round 35 --- 3 September 2026: validity is separated from
responsiveness in the four-dimensional cusp constructions}

\begin{itemize}[leftmargin=1.8em]
\item The native statement about one unmodified regular fibre is
\[
 (F_T,o_T)\longrightarrow
 (\mathbb T^2\times\mathbb R^2,(0,0)).
\]
It is derived directly from the original marked period family: two primitive
periods remain of unit length and the complementary rank-two quotient has
systole tending to infinity.

\item The later finite-cover theorem is valid but answers a different
question.  It changes the lattice and the base by
\[
 t_c=v^N,
 \qquad
 \Lambda^{(N)}=\mathbb Z\langle\widehat\gamma,\widehat u,
 N\widehat w,N\widehat\delta\rangle,
\]
and then lets both the cusp parameter and \(N\) diverge.  Its pointed
\(\mathbb R^4\) limit belongs to this altered covering family.  It neither
upgrades the original fibre limit nor uses the requested
\(S^{12}/S^{24}\) construction, so it is retained only as an ancillary
theorem.

\item The direct response to the multi-copy request is instead the cover-free
native quotient construction.  Two synchronized copies give
\[
 F_T^2\longrightarrow Q_T^2,
 \qquad
 (Q_T^2,(0,0))\longrightarrow(\mathbb R^4,0),
\]
and the based compactification globalizes its source through
\(X^{\wedge2}\cong S^{12}\).  Four synchronized copies of the unique
primitive full-monodromy-invariant circle quotient give the analogous
\(S^{24}\) route.  These statements quotient away compact directions; they
do not identify one original four-torus with \(\mathbb R^4\), and the named
spheres are domains of quotient maps rather than new compact complex
spheres.

\item The complements of the smash basepoints in the resulting moving
four-spheres are literally homeomorphic to \(\mathbb R^4\).  That fact alone
is the formal identity \(S^4\setminus\{*\}\cong\mathbb R^4\).  The
candidate-specific content is the exact factorization from the cusp quotient,
the placement of those moving four-spheres inside the fixed compact targets,
and the proof that their assigned intrinsic metrics exhaust Euclidean balls.
This distinction is now carried by the abstract, overview, specialist draft,
and provenance notes so that no nearby but easier theorem is substituted for
the requested construction.

\item The source transcript locates the exact missed globalization.  Its
two-copy calculation already contained the correct local map
\[
 F_T^2\xrightarrow{\ \pi_T\times\pi_T\ }Q_T^2
 \longrightarrow\mathbb R^4,
\]
but it stopped after proving the irrelevant obstruction
\(X\times X\not\simeq S^{12}\).  The global source used here is instead the
smash quotient \(X\wedge X\cong S^{12}\).  This is well typed because the
global based map \(q_2:(X,x_0)\to(\mathcal A,v_0)\), with \(x_0\) on the
central fibre, is constant on the global fat wedge.  The regular radial
subspace \(F_T^2\) is disjoint from that fat wedge, so
\(q_2^{\wedge2}\) restricts there to \(\pi_T\times\pi_T\) under the
product-to-smash embeddings.  By contrast, the fibrewise map based at the
identity does not factor through \(F_T\wedge F_T\); no such false
factorization is used.
\end{itemize}

\subsection*{Round 36 --- 3 September 2026: stable-stem nullity of both
sphere-valued carriers}

\begin{itemize}[leftmargin=1.8em]
\item The independently supplied based-homotopy calculation has been checked
against the cited primary locators.  Hatcher's suspension stable range
\cite[Corollary~4.24, p.~360]{Hatcher2002} gives
\[
 \pi_{12}(S^8)\cong\pi_4^{\mathrm S},
\]
and the \(k=4\) row of Isaksen--Wang--Xu's Table~1
\cite[p.~113]{IsaksenWangXu2023} is trivial, so
\[
 \pi_{12}(S^8)=0.
\]
The checked Isaksen--Wang--Xu source archive is
\texttt{C:/Users/user/Documents/math/literature\_primary/}%
\texttt{Isaksen\_Wang\_Xu\_2020\_stable\_homotopy\_source/}%
\texttt{arxiv\_2001.04511\_source.tar}; it has size \(108440\) bytes and
SHA-256
\texttt{178477986A1D07C6FAE12CA657BCA191D6B8903D98D78F3EB9650AEDC9AC25CB}.

\item For an arbitrary based map \(g:S^6\to S^3\), Hatcher's full
low-dimensional homotopy-group table \cite[p.~339]{Hatcher2002} gives
\[
 \pi_6(S^3)=\mathbb Z/12,\qquad
 \pi_{12}(S^6)=\mathbb Z/2.
\]
Thus \(12g=0\).  Toda's Proposition~5.11
\cite[Proposition~5.11]{Toda1962} identifies the latter group as
\[
 \pi_{12}(S^6)=\mathbb Z/2\{\nu_6\circ\nu_9\}.
\]
The same proposition also identifies
\(\pi_{n+6}(S^n)=\mathbb Z/2\{\nu_n\circ\nu_{n+3}\}\) for \(n\geq5\);
suspension takes each displayed generator to the next one.  Hence
\[
 \operatorname{st}:\pi_{12}(S^6)\xrightarrow{\ \cong\ }\pi_6^{\mathrm S}.
\]
Isaksen--Wang--Xu's Table~1 gives
\[
 \pi_3^{\mathrm S}=\mathbb Z/24\{\nu\},\qquad
 \pi_6^{\mathrm S}=\mathbb Z/2\{\nu^2\}.
\]
The relation \(12g=0\) forces
\(\operatorname{st}(g)=2a\nu\) for some \(a\in\mathbb Z/12\).
Multiplicativity of stabilization under smash products then gives
\[
 \operatorname{st}(g^{\wedge2})
 =\operatorname{st}(g)^2
 =4a^2\nu^2=0.
\]
The displayed stabilization isomorphism is injective, so
\([g^{\wedge2}]=0\in\pi_{12}(S^6)\).  Smashing this null-homotopy with
further copies of \(g\) proves
\[
 [g^{\wedge r}]=0\in\pi_{6r}(S^{3r})\qquad(r\geq2).
\]
Taking \(g=\widehat q_3\) proves that the full continuous surjective
carrier
\[
 \widehat q_3^{\wedge2}:S^{12}\twoheadrightarrow S^6
\]
is based null-homotopic.

\item Writing \(f=\widehat q_\gamma:S^6\to S^2\), the sixfold reduced
suspension
\[
 f\wedge\operatorname{id}_{S^6}=\Sigma^6f:S^{12}\longrightarrow S^8
\]
represents an element of the zero group above and is therefore based
null-homotopic.  The smash square has the exact factorization
\[
 f^{\wedge2}
 =\bigl(\operatorname{id}_{S^2}\wedge f\bigr)
  \circ\bigl(f\wedge\operatorname{id}_{S^6}\bigr):
 S^{12}\longrightarrow S^4.
\]
The right-hand factor is \(\Sigma^6f\), so \(f^{\wedge2}\) is based
null-homotopic.  For every \(r>2\), the canonical associativity
identification gives
\[
 f^{\wedge r}=f^{\wedge2}\wedge f^{\wedge(r-2)},
\]
and smashing the based null-homotopy of \(f^{\wedge2}\) with
\(f^{\wedge(r-2)}\) proves
\[
 [f^{\wedge r}]=0\in\pi_{6r}(S^{2r})
 \qquad(r\geq2).
\]
In particular, the already constructed continuous surjective quotient
\[
 \widehat q_\gamma^{\wedge4}:S^{24}\twoheadrightarrow S^8
\]
is based null-homotopic.  Its point-set quotient property, embedded moving
four-sphere, literal open \(\mathbb R^4\)-cell, and intrinsic metric limit
remain unchanged.

\item The scope boundary is explicit.  Neither calculation determines the
one-copy class \([\widehat q_3]\in\pi_6(S^3)\) or
\([\widehat q_\gamma]\in\pi_6(S^2)\).  The based homotopy classes of the maps into
\(\mathcal A\wedge\mathcal A\), \(\mathcal B^{\wedge4}\), and
\(\mathcal C^{\wedge4}\) also remain \Open.  The zero class of a
sphere-valued postcomposition does not imply that the preceding
conical-target map is null-homotopic.

\item The stable identifiers
\texttt{S6-CUSP-SPHERE-MOVING-S4},
\texttt{M-CUSP-SPHERE-Q3-SMASH},
\texttt{M-CUSP-SPHERE-QGAMMA-SMASH}, and
\texttt{U-CUSP-SPHERE-MOVING-S4} were updated in place, together with the
existing dependency row.  Consequently all machine-readable row counts and
status counts remain exactly those recorded at the end of Round~34: \(65\)
\texttt{VERIFIED}, \(4\) \texttt{OPEN}, and \(6\) \texttt{REFUTED}
claims; \(51\) \texttt{VERIFIED}, \(1\) \texttt{OPEN}, and \(2\)
\texttt{REFUTED} morphisms; \(31\) \texttt{VERIFIED}, \(5\)
\texttt{OPEN}, and \(4\) \texttt{REFUTED} uncertainties; and \(157\)
dependency rows.  No compilation was run during this propagation.
\end{itemize}

\subsection*{Round 37 --- 4 September 2026: exact Koide parametrization,
quaternion bridge, effective exponent, and \(10^8\) census}

\begin{itemize}
\item Integrated the complete HNF parametrization of ideals of norm
\(m_n=18n^2-2\), the unique seed \(a_0>2q_0\geq0\), the exact unit-orbit
inverse, the multiplicity formula
\(\rho(m_n)=\rho(3n-1)\rho(3n+1)\), the sharp floor formula, and the
Hecke-phase discrepancy splitting.
\item Verified effective well-roundedness for the literal chamber and height,
then applied Benoist--Oh, Theorem~13.2, to the finite integral-orbit union.
The resulting theorem is
\(R(X)=C_0X+O(X^{1-\eta})\) for an unspecified \(\eta>0\); no square-root
exponent is inferred.
\item Integrated the exact quaternion order, shell bijection, ramification,
centralizer, primitivity, finite conjugacy-class decomposition, and radial
Huber parameter.  Corrected the order description to reduced discriminant
\(108\), index \(18\) in a maximal overorder, and non-Eichler local type at
\(2,3\).  The principal congruence inclusion
\(\mathcal M^1(18\mathcal M)\subset\mathcal O^1\subset\mathcal M^1\)
provides the exact congruence-lattice bridge.
\item Integrated the independently replayed exact census through \(10^8\):
\(R(10^8)=3{,}536{,}002\), \(2{,}327{,}513\) positive-\(n\) canonical
unit orbits, and all
\(3{,}536{,}002\) coordinate rows rechecked.  These finite data are not used
as an asymptotic proof.
\item Preserved the status of
\(R(X)-C_0X=O_\varepsilon(X^{1/2+\varepsilon})\) as \Open.  The
primary-source audit separates the exact unspecified power saving, the smooth
radial \(39/64\) estimate, the sharp radial \(2/3\) estimate, and the missing
sharp nonradial boundary/complementary-spectrum input.
\item Added the distinct stable claim
\texttt{S6-KOIDE-CENSUS-1E8} with status \Verified.  Its object is exactly
\[
 a^2-3q^2-18n^2=-2,\qquad a>0,\quad q,n\geq0,\quad a+4n>q,
 \qquad 4a+17n\leq X,
\]
without a normalized or substituted height.  It records
\(R(10^8)=3{,}536{,}002\), positive-\(n\) endpoint \(2{,}943{,}725\),
\(2{,}327{,}513\) positive-\(n\) canonical unit orbits plus the separate
\(n=0\) orbit, and a full
independent recheck of every emitted six-coordinate row.  The dependency
ledger now records both its use of \texttt{S6-KOIDE-DEFECT-COUNTING} and the
fact that it does not imply \texttt{S6-OPEN-KOIDE-SQRT-DISCREPANCY}.
\item Promoted and checked the project-relative reproducibility packet
\texttt{supporting\_materials/workbench/evidence/koide\_exact\_checkpoint\_20260904/}.
It contains \(18\) files.  The self-excluding \(17\)-entry manifest has SHA-256
\texttt{8D108EB928E5D07093BADDE169F443A4A38B6FCF3EB59D5192BDD33ECEF8D7A4}
and verifies with zero mismatches; its README has SHA-256
\texttt{FFF150860F000CF17DCD0CD3928D939BFA5B501D59C715DE444A3A8D8CDF899A}.
The final census result JSON, compressed full coordinate stream, and
coordinate-audit receipt have SHA-256, respectively,
\texttt{DA92F0866074E530B8174A48915CAA0F65CD4CD8957E209306ADC200739D1961},
\texttt{DFC313492AEDB396DCFA10520BFEFC699BA601357BC9933D921DFA3884310BD4},
and
\texttt{3010363108E7B460B7331A92F3263BBD77992AEDFC3064973750E16D707DA80F}.
\item Replayed the bundled parametrization and quaternion verifiers.  Both
exited with status zero: the former checked the HNF count, shifted character
factorization, unique seeds, inverse descent, and brute-force counts on its
declared bounded range; the latter returned \texttt{PASS}, checked \(48\)
finite-box shell points, determinant \(-11664\), reduced discriminant \(108\),
and the Hilbert symbols at \(2\) and \(3\).  These bounded checks do not
replace the canonical ideal, centralizer, conjugacy-class, or spectral proofs.
\item Refreshed the current Koide source provenance after the final precision
repair.  The SHA-256 values are
\texttt{965B0A3DCA8BFE4F22B9E9208368FA04697070FD25BD06BD67F2481D17CB9D9D}
for \texttt{koide\_counting\_repair.tex},
\texttt{A457709DCCE0FCE0774EB9326B84B00D5D8144C2087F480DA704974E499BCDD3}
for \texttt{koide\_exact\_parametrization.tex},
\texttt{18ED5968B7E67D9A0B23D8675B40C37BA46170AD29EFDBB9CE10ED632C99600A}
for \texttt{koide\_exact\_census.tex},
\texttt{93759F7BA99870C84C9144A373A53D2E9A8BC3AE909518E47991373ACE0B47BB}
for \texttt{koide\_quaternion\_bridge.tex}, and
\texttt{BC65B7574FD51DA36EE8B5B594720BD357187EA16AFEFA62F07110C3CEBB333E}
for \texttt{koide\_effective\_literature\_audit.tex}.
\item After this ledger update, the machine-readable totals are \(81\)
claims (\(71\) \texttt{VERIFIED}, \(4\) \texttt{OPEN}, \(6\)
\texttt{REFUTED}), \(60\) morphisms (\(57\) \texttt{VERIFIED}, \(1\)
\texttt{OPEN}, \(2\) \texttt{REFUTED}), \(40\) uncertainties (\(31\)
\texttt{VERIFIED}, \(5\) \texttt{OPEN}, \(4\) \texttt{REFUTED}), and
\(168\) dependency rows.  JSONL parsing, identifier uniqueness, dependency
target resolution, evidence-path existence, and cited checkpoint hashes all
passed.  No research module, root TeX file, metadata file, or PDF was edited
by this ledger pass, and no compilation was run.
\end{itemize}

\subsection*{Round 38 --- 4 September 2026: independent audit of the exact
Lorentz-observation bridge}

\begin{itemize}
\item Promoted the complete independent audit into
\texttt{supporting\_materials/workbench/evidence/
exact\_lorentz\_observation\_bridge\_audit\_20260904/}.  Its
\texttt{PROMOTION\_SHA256SUMS.txt} has SHA-256
\texttt{52F53746558FA219260B224ABDB71C86D2CD6C13192774D0D2FB1A1ED71BCF35}
and verifies all \(14\) non-self manifest rows.  The packet TeX is byte-for-byte
identical to the already integrated canonical evidence source: \(53{,}061\)
bytes, SHA-256
\texttt{2BE16E1DC71B2E5DC04238F0FBE440138258D0825A4841AB2F29662587681599}.

\item Replayed the packet checker with all \(33/33\) checks passing and ran a
separate verifier with all \(30/30\) checks passing.  The resulting
claim-by-claim audit report has SHA-256
\texttt{C8D01AF28836340798C2FA7E27F8E097D67E0B8D7C13F75271AF473958D391DF}.
It proves packet claims C01--C22, C24--C26, and C31, as well as the
all-spin occurrence part of C30.  The proposed unique global spin selection
in C30 remains \Open, and C27--C29 remain \Open.

\item The audit confirms the exact rank stratification of the hard-constrained
momentum map without replacing its fixed parameters or its strata.  For
\(w\neq0\), \(Dp\) has rank \(4\), its kernel is the tangent space of the
three-sphere momentum fibre, and the intrinsic action gives
\(\dot p=0\) and \(\dot x=0\).  At \(w=0\), \(Dp\) has rank \(3\), its
image is \(p^{\perp_L}\), the momentum fibre is a point, and the equations are
\(\dot p=0\) and \(\dot x\in\mathbb Rp\).  This refutes C23 when read as
propagating worldline dynamics on every stratum and C32 when read as a global
completion theorem.

\item The stable records
\texttt{S6-WPR-HARD-THRESHOLD-DYNAMICS} and
\texttt{M-WPR-HARD-PACKET-MOMENTUM} now cite both independent checker outputs,
the claim audit, the promotion manifest, and the hash-pinned source.  The
separate multiplication/Fock spectrum still does not define the Hamiltonian
of the intrinsic worldline action, so neither the determinant threshold nor
the free spectral calculation proves a nonlinear interacting Yang--Mills
mass gap.
\end{itemize}

\subsection*{4 September 2026: complete divisor-channel intake and rational continuation}
The immutable Koide--Erd\H{o}s--Straus v3 archive was verified against its
published hash and preserved in
\path{supporting_materials/workbench/evidence/koide_es_v3_audit_20260904/source/}.
Its original verifier was replayed in a separate directory. An independent
smaller-denominator enumeration agrees on 1430 ordered solutions in 665
coprime shells; the diagonal count gives 819 certificates for 329 primes.
The exact proofs, the positive-real-order counterexample, and an additional
rational inverse-bearing map to the defect-one Koide surface are in
\path{supporting_materials/workbench/research/koide_es_v3_audit.tex}.
For integral source entries the new map has denominator
$|D|/\gcd(|D|,|2K|)$, with
$D=-1-a^2-b^2+4a+4b+4ab$, $K=1-2a-2b$.
The earlier orthogonal map is retained and still has no rational image.
Elsholtz--Tao's primary LaTeX (arXiv:1107.1010v6, Appendix A) was obtained
and read for the classical shifted-prime density context. No global
threefold, all-prime, or physical conclusion is inferred from these maps.
The six-page standalone audit was compiled and all six rendered pages
visually inspected. The source and feedback are mirrored to the existing
canonical Overleaf; no new project or sharing change is involved.

\subsection*{4 September 2026: exact reconciliation of the quaternionic continuations}
The two supplied Lorentz thirteen-slices are proved to coincide as embedded
subspaces of $\operatorname{Herm}_3(\mathbb H)$, through
$PB(b)P^{\mathsf T}=\overline{u^*b}\varepsilon$ and
$X^{-1}B(b)X^{-\mathsf T}=B(X^*b)$. The second-spinor dictionary is
$w=-\varepsilon\bar\chi$, with inverse $\chi=\varepsilon\bar w$.
The complete cubic, positive fibres, sphere pair $(S^{12},S^8)$, and
fifteen-dimensional Jordan closure are proved in
\path{supporting_materials/workbench/research/quaternionic_4441_reconciliation.tex}.
The Jordan product retains $D=X^*X=\operatorname{diag}(C^{-1},3)$ and the
explicit nonclosure witness has coefficient $3/2$, not $1/2$.
For the actual supplied scalar kinetic terms and potential, initial data
$\Phi=\operatorname{diag}(1,-1,0)$, $b=(1,-1,0)^{\mathsf T}$, zero velocities,
give $\rho=\dot\rho=0$ but $\ddot\rho=-2\kappa$.
The full model and the different constrained variational model are both
retained, with reaction $\kappa(uu^*/3)\Phi b$. This is not an assertion
that the original packet confused symmetry invariance with all dynamics;
it makes their distinction exact. The ordinary left group action uses
$gZk^{-1}$, and anti-Hermitian gauge generators require the positive trace
coefficient for the displayed $(+---)$ scalar convention.
The evidence archive preserves 24 original files, verifies 16 source
manifest rows, and records restoration of four replay receipts whose line
endings had changed. The 26-check independent script uses its own
quaternionic matrix multiplication. Wang's primary source and the local
Alesker text were checked for the spectral and determinant conventions.
The broader vacuum and quantization claims remain outside this bounded audit.

\subsection*{4 September 2026: global invariant scalar sectors and fourth-order escape}
The original scalar equations preserve the fixed-subspace configuration
spaces $\mathcal T_W=\{(\Phi,b):[\Phi,E_W]=0,b\in W,\operatorname{tr}\Phi=0\}$,
for complex $W\subset u^\perp$ of dimensions one or two, with tangent
velocities. Their dimensions are six and eight; their reactions vanish.
Energy conservation and an explicit quartic coercivity estimate prove
global existence. The line example retains $s_0=42$, $\chi_0=(1,0)$ and
$z_0=(0,3)$, so the construction is not restricted to $z=0$.
The exact cancellation in $c=u^*\Phi b$ gives
$\ddot c=-2(m^2+\Gamma s)c-(\kappa s/3)\rho+2u^*\dot\Phi\dot b$.
A separate exact witness has vanishing normal derivatives through order
three but $\rho^{(4)}(0)=-12\kappa$. A bounded independent agent checked
the block proof and normal jets, and a separate root verifier passes
20 exact checks. The maximal invariant set is not classified by these
results; the next Pro request concerns possible moving-subspace solutions.
\subsection*{4 September 2026: periodic moving kernels, colour moment, and scalar vacua}
The fixed-$W$ invariant families are not exhaustive. An explicit periodic
solution at $(m^2,\Gamma,\kappa)=(1,1,6)$ has $s=1$, $\rho=0$ and
$\Phi b=0$, but its rotating $b$ spans $u^\perp$, which $\Phi_0$ does
not preserve. A continuous parameter family is proved in the original
coordinates. The conserved moment is $\mathcal J=-s\Omega/2$; the
concrete orbit has energy $7/4$ and $\mathcal J=-\Omega/2\ne0$.
Variation of the anti-Hermitian $A_0$ shows why its zero-connection lift
fails Gauss' equation. The exact moving-frame transformation makes scalar
coordinates static while retaining $A'_0=\Omega$ and the same residual.
This obstruction is to that lift, not to every possible gauge-field lift.
The triality packet's scalar global-minimum theorem and full metric Hessian
are independently checked and proved in the canonical supplement, including
the sharp discriminant identity, all equality cases, coexistence, endpoints,
and all fourteen Hessian directions. Gauge quotients and finite-lattice
quantization are not promoted by that scalar calculation. Tong's primary
Gauss-law derivation was read for the gauge comparison.
\subsection*{4 September 2026: full-carrier observation and actual period comparison}
The full positive quaternionic carrier has a split positive observation to
$\operatorname{Herm}_2(\mathbb C)$ whose exact permitted-pair image is
$G\succeq0$, $|\rho|^2\leq\det G$. All permitted pairs lift, and
explicit positive lifts approach determinant zero with $\rho\ne0$.
Trace-free scalar matrices are not positive except at zero; their exact
spectral-floor observation retains $PP^*$. The rank-six local toroidal
period extension also gives a commuting positive Gram square. The
constant-metric obstruction and a descending metric for arbitrary positive
$W_*$ are proved, including every cusp-correction term. The resulting
leading Lorentz determinant is strictly positive, not null. These
constructions identify exact maps across the dimension experiments without
identifying their sphere boundaries or claiming a physical mass spectrum.

\paragraph{4 September: full field equations and retained-period metric fibres.}
Derived every homogeneous gauge--scalar equation before gauge fixing.
The literal periodic orbit has no temporal-gauge lift with arbitrary
homogeneous spatial matrices; its residual is a positive coefficient
times $\Omega$. Arbitrary temporal connection is retained in the exact
Gauss-inertia operator; pointwise solvability is not promoted to a full
field solution. A real radial sector solves all field equations globally.
The colour-group average of the Lorentz determinant is
$V_\mu\operatorname{tr}\operatorname{adj}(A)$, equal to $V_\mu\|b\|^2$
on the positive rank-one boundary; explicit gauge-inequivalent minima
vary this invariant at fixed potential. Separately, the unchanged period
map has complete positive lattice-cometric/weight fibres of dimensions twelve and seventeen.
An explicit rank-losing metric retains all cusp constants and has
condition number asymptotic to a displayed coefficient times $y^{p+2}$.
Bounded relative condition number obstructs that loss. Primary source
context checked against Tong and Gambioli; exact checkers have 89 and
61 assertions. See the cumulative quaternionic supplement for all proofs.

\paragraph{4 September: integrated full-gauge continuations and native Hermitian spectra.}
The exact invertible map
$C=K^*O^*\mathcal A_0OK+A$, $Y_i=K^*O^*\mathcal A_iOK$
retains the literal scalar fields as constant $H,q$ and preserves every
field equation (Proposition~\ref{q444:static-lift-system}). The full scalar
equations force, for real-skew connections, $\dot c=0$ and
$Q=cc^{\mathsf T}-\sum_i y_iy_i^{\mathsf T}$ with $Qe_1=0$ and
$\operatorname{tr}Q=a>0$. Covariance positivity forces $c,y_i\perp e_1$;
the exact Gauss pairing
$2(r^2+B^2)|c|^2+(4/g^2)\sum_i|y_i\times c|^2=0$
then gives a contradiction. This excludes arbitrary real spatial time
dependence, not just temporal gauge. Constant $Y_i$ in the full complex
algebra are separately excluded for arbitrary $C(t)$ by positive Gauss
inertia and $\sum_i|Y_iq|^2+aB^2=0$. Both exclusions extend through
the stated matter-stabilizer gauge transformations; arbitrary complex
time-dependent fields outside those classes remain open. Complete proofs
are Theorem~\ref{q444:real-no-lift} and
Corollaries~\ref{q444:static-complex-no-lift}--\ref{q444:real-stabilizer-no-lift}.

The retained vacuum tangent has squared norm $18\chi^2$ and nonzero
normal acceleration, so a fixed-amplitude homogeneous vacuum curve is not
an unconstrained solution (Proposition~\ref{q444:vk-horizontal}). The
independent exact sector
$\Phi=\operatorname{diag}(h+k,h-k,-2h)$, $b=\beta e_3$,
$\mathcal A_\mu=0$ satisfies every field equation with all currents zero.
Theorem~\ref{q444:vk-periodic} gives finite-amplitude periodic solutions
with $E-V_0=9h_0^2\omega^2+O(\omega^4)$ as $\omega\to0$;
Proposition~\ref{q444:vk-wave} gives exact nonlinear null waves
$q_{f(t-x_1)}$ and their retained-volume energy. Theorem~\ref{q444:gauge-spectrum}
computes all eight inertia eigenvalues with the actual generator metric,
and verifies transverse modes against all linearized equations including
Gauss. The real conserved invariant current
$j_\nu=-i(b^*D_\nu b-(D_\nu b)^*b)$ detects the diagonal branch with
$\mu_{\mathrm d}^2=(2g^2/3)\|b\|^2$, equivalently the retained-volume
Haar determinant formula in Proposition~\ref{q444:current-mass}.
These are classical statements in the specified model, not claims about
the interacting quantum spectrum.

The positive period weight $W$ is now identified exactly as the lattice
cometric, not the covariant lattice metric; that metric is $W^{-1}$.
Hermitian compatibility removes the twelve anti-complex fibre directions
and gives $h_3=2H^{-1}$ (Theorem~\ref{q444:hermitian-cometric}).
The original integral torus quotient has metric
$h_2=2(PHP^*)^{-1}$ and an exact Laplace intertwiner
(Theorem~\ref{q444:metric-submersion}). The unchanged lattice has complete
marked spectrum $4\pi^2k^{\mathsf T}Wk$, which recovers $W$
(Theorem~\ref{q444:torus-laplace}). Native $h_3=I_3$ gives $H=2I_3$,
$G=2PP^*$, and $h_2=(PP^*)^{-1}$; the source and target first positive
gaps are respectively asymptotic to $4\pi^2/y^2$ and $8\pi^2/y^2$,
with every period correction and $y=-\log|t|/(2\pi)$ retained
(Proposition~\ref{q444:native-period-spectrum}). No lattice, metric scale,
global transition, or quantum Hamiltonian is substituted by these maps.

\subsection*{4 September: quantum confinement of all homogeneous fields}
Section~\ref{qh:chapter} derives the full $24+8+6$ coordinate
Hamiltonian, its Friedrichs domain, all eight Gauss generators and
their common-kernel Hilbert space. The scalar estimate
$|\det\Phi-b^*\Phi b|\leq s^{3/2}/(3\sqrt6)$ is proved directly
from the complex-Hermitian eigenvalues without imposing a minimum
curve. Six directed oscillator slices give
\eqref{qh:full-form-bound}; a finite-cell approximation proves
compactness of the complete coupled form domain. The strictly
positive Brownian-bridge kernel makes the ground state simple and
invariant, and hence gives the strictly positive finite-model
physical gap in Corollary~\ref{qh:hq-ground}.
No numerical gap or volume-uniform estimate is asserted.

Proposition~\ref{qh:hq-embedding} retains the action and symplectic
pullbacks into torus fields. Its orbit-space map is not injective:
a periodic spatial gauge transformation maps the zero connection
to a nonzero homogeneous flat connection. The mixed Fourier
potential $3\Gamma\mathcal V_3q^2y^2/4$ also obstructs the proposed
all-homogeneous-wavefunctions tensor fixed-mode-vacuum quantum
subsector. The proof includes its operator-domain argument.
These calculations specify the next continuum comparison instead
of inferring it from classical low frequencies.

\subsection*{4 September: unrestricted rotating equations and a moving quantum comparison}
Theorem~\ref{q444:cx:thm:block} excludes a genuinely complex
cross--cross--block rotating pair with every phase and its $SO(2)$
spatial saturation retained. The full inverse in
Theorem~\ref{q444:complex-exact-zero-set} removes $C,\omega$ without
discarding the full scalar vector residual or either spatial residual.
The complete scalar-coordinate equivalence and quantitative stabilizer
constraint (\ref{q444:complex-full-scalar-equivalence},
\ref{q444:complex-stabilizer-spacelike}) apply outside that class.
The stress equation in Proposition~\ref{q444:complex-stress-obstruction}
adds a positive matrix and sum-of-squares test. These are exact
constraints, not a solution or emptiness claim for the reduced zero set.

Section~\ref{fq:section} replaces the failed fixed-mode vacuum map by
the actual positive quartic ground-state line. Its exact form pullback
retains the Born--Huang potential, and its full off-diagonal residual
is proved on the effective operator domain. A min--max bound proves
the needed fibre separation in \eqref{fq:regime}; bounded multipliers
give actual quartic moments rather than Gaussian substitutions.
The explicit physical-time error \eqref{fq:time-error} follows by
Duhamel with all original volume and Planck factors. It is not a
higher-order adiabatic claim or an estimate uniform as $\omega_*\to0$.

\subsection*{5 September: the original cusp coefficient is zero}
The complete framing calculation in Theorem~\ref{cusp-null:thm}
(also Theorem~\ref{h24:thm:cusp-g-zero} in the independent higher-rung
draft) places the entire regular inverse image in one original toric
chart. The reciprocal-unit coordinate map retains all of $C(t)$ and
has the explicitly written inverse. The actual dual frame is $NJH^{-1}$;
its invertible homotopy extends across the collared $D^2\times T^2$.
The relative smoothing argument uses a single fixed closed neighborhood
for the positive separation bound. Thus $g\simeq *$ before sphere
recognition, so $a=0$, $E^6[g]=0$, and $\Theta_g\chi=0$.
The universal suspension $\mathbb Z/12\hookrightarrow\mathbb Z/24$
remains the even-subgroup embedding; the particular map represents its
zero element. The explicit construction does not assert nullity of
the earlier conical target or of the distinct one-copy map to $S^2$.

\subsection*{5 September: exact finite coupling fibres and generic spatial exclusion}
Theorem~\ref{q444:finite-coupling-fibre} gives the entire coupling
fibre over each unrestricted spatial pair, with at most fourteen
positive values and the full both-sign inverse.
Proposition~\ref{q444:finite-coupling-exact-degree} proves degree
exactly twenty-eight whenever the original domain polynomial is
nonzero, by a positive leading-coefficient formula.
Theorem~\ref{q444:generic-projection-obstruction} constructs a nonzero
polynomial whose zero set contains the entire spatial projection.
Its exact witness is checked by rational Bezout arithmetic and an
independent full-eight-dimensional finite-field reconstruction.
The original 31-check and independent 45-check suites pass, as do
all 101 modular evaluations and the independent coefficient identities.
These finite certificates supplement the complete universal proofs.
The algebraic exceptional set has not been proved empty or to contain
a solution. No arbitrary-time gauge or continuum quantum conclusion
is inferred. Original candidate and reviewer provenance are unchanged.

\subsection*{5 September: a full exceptional-class exclusion and quantitative certificate}
The complete two-cross proof is included in
Section~\ref{q444:twocross:section}. The linear spatial class
$\mathfrak p\times\mathfrak p$ lies inside the exceptional zero set,
yet admits no full solution at positive original coupling.
Gauss forces the temporal block, the scalar equations classify it with
two exact inverse frame maps, and every gauge residual is calculated.
The final positive contradiction retains both signs of the frequency.
Proposition~\ref{q444:twocross:quantitative} strengthens this to a
denominator-free polynomial residual certificate and a uniform spatial
error bound on bounded scalar--Gauss loci with $g^2$ bounded away from zero.
Both proofs received independent equation-by-equation review. The original
26 exact checks and two additional certificate checks passed. Existence
on the remaining exceptional set, or of more general time dependence,
is not settled by these results. Historical candidate annotations are unchanged.

\subsection*{5 September: direct-image and logarithmic real-marking corrections}
The independent bounded recheck found two errors in the written proof,
both corrected in the current source. Theorem~\ref{thm:cartier-duality}
now constructs the pairing through the analytic Gorenstein dualizing
complex and trace. Matching outer terms of exact sequences is not a
construction of a middle duality morphism. The complete replacement
retains the canonical twist and the cohomological shift.

The holomorphic overlap translation by $\alpha\Pi(s)v_j$ has real
vector $a_j=\operatorname{Re}\alpha\,v_j+\operatorname{Im}\alpha\,J_sv_j$.
Substitution includes both the mixed term and a horizontal two-form.
On a full clockwise covering circle the latter vanishes and
$\Delta a_j=-v_j$, giving $N_{T_j}d_j=\iota_{v_j}E$.
Period covariance gives $a_j(g_js)=A_ja_j(s)-v_j/m_j$, hence the
same based affine endpoint and meridian power relation.
These replace the incorrect complex-scalar expression in a real marking.
The conductor-to-fibre-cohomology-to-direct-image implication survives
this audit; a global counterexample to CDP20 is not inferred merely
from formal proof defects. Exact before/after sources, primary-source
locators and the audit are retained in
\path{supporting_materials/workbench/evidence/s6_cdp_direct_image_reaudit_20260905/}.

\subsection*{5 September: complete scalar--Gauss domain and original coupling surjection}
The complete proof in Section~\ref{q444:sgd:section} classifies the entire
two-cross scalar--Gauss zero set. With $\delta=\sqrt{2/3}$,
\[
 \mathcal D=\{(u,v_0,y):u>0,\ d=uv_0-\delta^2-y^2>0,
 c=4v_0-5u>0\},\qquad t=72d/c,
\]
Theorem~\ref{sgd:parametrization} proves the classified bijection
$\mathcal D\times\{1,-1\}^2\times U(1)^2$ with explicit inverse signs
and phases. Lemma~\ref{sgd:literal-addon:temporal} derives the arbitrary
literal temporal form from Gauss and scalars alone. The full literal map
of Theorem~\ref{sgd:literal-addon:quotient} adds $\zeta\in U(1)$ and has
exactly the two-element fibres of
$(p_+,p_-,\zeta)\mapsto(-p_+,-p_-,-\zeta)$.
For each prescribed original $\tau=g^2>0$,
\[
 \rho=\frac{\tau+\sqrt{\tau^2+4608\delta^2}}{96\delta},\qquad
 M_\rho=\delta\begin{pmatrix}\rho&1\\1&2\rho\end{pmatrix},\qquad
 48\delta\rho^2-\tau\rho-24\delta=0
\]
prove $t=\tau$, with both temporal signs, both original frequency signs,
and every scalar and Gauss entry zero (Corollary~\ref{sgd:every-coupling}).
This supplies a nonempty domain for the previously proved quantitative
spatial-error certificate. The two spatial residuals cannot vanish simultaneously
throughout this class, and unrestricted exceptional-set existence remains
open. No quantum gap conclusion follows.

The independent review passed the classified inverse, the full literal
quotient proof, and all 25 exact replay groups. Candidate, addendum and
review are preserved in
\path{supporting_materials/workbench/evidence/scalar_gauss_domain_20260905/}.
The old statement that the residual certificate did not assert
nonemptiness is superseded only on current surfaces; historical annotations
and original proof snapshots retain their provenance. This entry records
source integration; serial builds and publication have separate receipts.

\subsection*{5 September: the residue lattice, full Weyl quotient, and exact shell maps}
Section~\ref{esn:sec:shared} constructs
$R_{\mathrm{ES}}=\bigoplus_{\eta\in\mathbb F_2^2}A(F_\eta)\perp D_4$
in the original six-coordinate fibre order. The full discriminant,
order-$72$ isotropic glue, and all glue-class minima prove
$[N:R_{\mathrm{ES}}]=72$, $\det N=1$, and exactly the original
$144$ roots of type $A_5^4D_4$
(Theorems~\ref{esn:thm:root-discriminant} and~\ref{esn:thm:niemeier}).
The coordinate transporter preserves the earlier glue markings. The
$D_4$ shortest dual classes $v,s,c$ match the vector, even half-spin,
and odd half-spin weight sets of the flag tangent modules described
by Mare--Willems \cite{esn:MareWillems2013}.

Theorems~\ref{esn:fa:thm:full-quotient} and~\ref{esn:fa:thm:residue-weyl}
prove the exhaustive factorization
$O(N)=W(R_{\mathrm{ES}})\rtimes L(G)$, $G\cong GL_2(\mathbb F_3)$,
and $\rho(S_{840})\subset W(R_{\mathrm{ES}})$ with
$\rho(r)L_g=L_g\rho(r)$ for this chosen complement.
The $D_4$ discriminant map $G\to GL_2(\mathbb F_2)$ has kernel $Q_8$
and the displayed $S_3$ splitting. Cheng--Duncan--Harvey already record
the abstract Weyl quotient \cite[Table~3]{esn:CDH2014}; the source here
supplies the exact code isomorphism, coordinate lift and full factorization.

For $p\equiv1\pmod{24}$ prime, $R>0$, $R\equiv3\pmod4$, retain
$a=(p+R)/4$, $g=\gcd(R,pa)$, $R_0=R/g$ and $N_0=pa/g$.
Theorem~\ref{esn:thm:shell-bijection} proves the bijection from
$\{d>0:d\mid N_0^2,\ d\equiv-N_0\pmod{R_0}\}$ to all positive
ordered pairs solving $4/p=1/a+1/b+1/c$, by
$(b,c)=((N_0+d)/R_0,(N_0+N_0^2/d)/R_0)$, with inverse $d=R_0b-N_0$.
Its unordered count is $(A_R(p)+\mathbf1_{R_0\mid2})/2$.
At the fixed shell $R=3$, the residue-one primes $2521,3361$ have
raw counts $0,6$ and unordered counts $0,3$. For a finite prime sample
$P$ reducing into $S_{840}$, Theorem~\ref{esn:thm:sample-fourier}
proves Fourier inversion for the aggregate function $w_{R,P}$ and the
integral map $Q_{R,P}=6W_{R,P}-M_{R,P}u\in A(S_{840})\subset N$.
The actual shell and sample remain in the domain. The complete proofs
are included once in each of the workbench, ES and higher-rung hosts,
with one shared bibliography helper and accompanying source provenance.
This entry records source integration; builds and publication have separate
receipts. These results do not settle universal selector nonemptiness,
the global higher torus wrap, or a physical interpretation.
